%%%%%%%%%%%%%%%%%%%%%%%%%%%%%%%%%%%%%%%%%%%%%%%%%%%%%%%%%%%%%%%%%%%%%%
% Lớp 11 | Hàm số mũ và hàm số lôgarít | Phương trình, bất phương trình mũ và lôgarít
% Dạng 1D6?4-6: Toán thực tế áp dụng phương trình mũ, lôgarít
% Số câu: 102  (sinh tự động bởi tach_cau.py; từ nay thầy sửa trực tiếp file này)
%%%%%%%%%%%%%%%%%%%%%%%%%%%%%%%%%%%%%%%%%%%%%%%%%%%%%%%%%%%%%%%%%%%%%%

%%%%%%%%%%%%%%%% TRẮC NGHIỆM 4 PHƯƠNG ÁN %%%%%%%%%%%%%%%%

% nguồn: Ok_Toan11_TracNghiem_Lan2
\begin{ex}%[1D6H4-6]%[ma:e4d092a2]
	Bác An gửi tiết kiệm số tiền $300$ triệu đồng kì hạn $12$ tháng với lãi suất kép $5\%$ một năm. Giả sử lãi suất không thay đổi. Sau ba năm vì cần tiền nên bác An đến ngân hàng rút ra $100$ triệu đồng, phần còn lại vẫn tiếp tục  gửi. Hết bốn năm tiếp theo, bác An lại đến ngân hàng rút toàn bộ tiền tiết kiệm (cả gốc và lãi) về, hỏi bác An sẽ thu về được bao nhiêu tiền (đơn vị: triệu đồng, kết quả làm tròn đến hàng phần nghìn)?
	 \choice
		 {$294{,}000$}
		 {$296{,}688$}
		 {\True $300{,}580$}
		{$225{,}178$} 
	 \loigiai{
		Số tiền sau $3$ năm là 
		 \[
		A_1=P_1(1+r)^{N_1}=300(1{,}05)^3. 
		 \]
		Bác An rút $100$ triệu đồng, số  tiền còn lại là
		 \[
		 P_2=A_1-100.
		\] 
		Sau $4$ năm tiếp theo số tiền cả  gốc và lãi bác An thu được là
		\[ 
		A_2=P_2(1{,}05)^4\approx300{,}580\text{  (triệu đồng).}
		\] 
	 }
\end{ex}

% nguồn: Ok_Toan11_TracNghiem_Lan2
\begin{ex}%[1D6H4-6]%[ma:a157b8b6]
	Số người trong cộng đồng sinh viên  đã nghe một tin đồn nào đó là $N=P\left(1-\mathrm{e}^{-0{,}15d}\right)$, trong đó $P$ là tổng số sinh viên của cộng đồng và $d$ là số ngày trôi qua kể từ khi tin đồn bắt đầu. Trong một cộng đồng $1\,000$ sinh viên, cần bao nhiêu ngày để $450$ sinh viên nghe được tin đồn?
	\choice 
		 {\True $4$}
		{$3$} 
		 {$5$}
		{$2$} 
	 \loigiai{
		 Ta có
		\begin{eqnarray*}
			&&N=P(1-\mathrm{e}^{-0{,}15d})\\ 
			&\Leftrightarrow&450=1\,000(1-\mathrm{e}^{-0{,}15d})\\ 
			&\Leftrightarrow&\mathrm{e}^{-0{,}15d}=\dfrac{11}{20}\\ 
			 &\Leftrightarrow&-0{,}15d=\ln\left(\dfrac{11}{20}\right)\\
			&\Leftrightarrow&d=\dfrac{\ln\left(\dfrac{20}{11}\right)}{0{,}15}\approx3{,}98. 
		\end{eqnarray*}
		Vậy cần $4$ ngày để $450$ sinh viên  nghe được tin đồn.
	 }
\end{ex}

% nguồn: Ok_Toan11_TracNghiem_Lan2
\begin{ex}%[1D6H4-6]%[ma:fd5b2672]
	Biết rằng năm $2001$, dân số Việt Nam là $78\,685\,800$ người và tỉ lệ tăng dân số năm đó là $1{,}7\%$. Cho biết sự  tăng dân số được ước tính theo công thức $S=A\mathrm{e}^{Nr}$ (trong đó $A$ là dân số của năm lấy làm mốc tính, $S$ là dân số sau $N$ năm, $r$ là tỉ lệ tăng dân số hằng năm). Cứ tăng dân số với tỉ lệ như vậy thì đến năm nào dân số nước ta ở mức $120$ triệu người?
	\choice 
		{$2029$} 
		 {$2020$}
		{$2025$} 
		 {\True $2026$}
	\loigiai{ 
		Ta  có
		\begin{eqnarray*}
			 &&S=A\mathrm{e}^{Nr}\\
			&\Leftrightarrow&N=\dfrac1r\ln\dfrac{S}{A}. 
		\end{eqnarray*}
		Với $A=78\,685\,800$, $r=0{,}017$, $S=120\,000\,000$, ta  được
		 \[N=\dfrac1{0{,}017}\ln\dfrac{120\,000\,000}{78\,685\,800}
		 \approx24{,}83.\]
		Vậy sau $25$ năm, tức là vào năm $2001+25=2026$, dân số nước ta đạt mức $120$ triệu người. 
	 }
\end{ex}

% nguồn: Ok_Toan11_TracNghiem_Lan2
\begin{ex}%[1D6H4-6]%[ma:34b8aa5e]
	Bác An gửi tiết kiệm số tiền $100$ triệu đồng kì hạn $12$  tháng với lãi suất kép $5\%$ một năm. Giả sử lãi suất không thay đổi, tính số tiền (cả vốn lẫn lãi) bác An thu được sau $5$ năm (đơn vị: triệu đồng, kết quả làm tròn đến hàng phần nghìn).
	\choice 
		 {$125{,}628$}
		{$130{,}432$} 
		{\True  $127{,}628$}
		{$125{,}000$} 
	\loigiai{ 
		Áp dụng công thức lãi kép $A=P(1+r)^N$ với  $P=100$, $r=0{,}05$, $N=5$, ta có
		 \[
		A=100(1{,}05)^5\approx127{,}628\text{ (triệu  đồng).}
		\] 
	 }
\end{ex}

% nguồn: Ok_Toan11_TracNghiem_Lan2
\begin{ex}%[1D6H4-6]%[ma:980f639c]
	Tại thời điểm ban đầu nếu đầu tư $P$ đô la với tỷ lệ lãi suất được tính gộp liên tục hằng năm không  đổi là $r$ thì giá trị tương lai của khoản đầu tư này sau $t$ năm là $B(t)=P\cdot \mathrm{e}^{rt}$ đô la. Giả sử tỷ lệ lãi suất tính gộp hằng năm là $8\%$. Hỏi sau bao nhiêu năm thì số tiền đầu tư ban đầu tăng thêm ít nhất $50\%$.
	 \choice
		{$5$} 
		 {$8$}
		{$7$} 
		 {\True $6$}
	 \loigiai{
		Ta có 
		\[P\cdot \mathrm{e}^{0{,}08t}>1{,}5P\Leftrightarrow  \mathrm{e}^{0{,}08t}>1{,}5\Leftrightarrow 0{,}08t>\ln1{,}5\Rightarrow t>\dfrac{\ln1{,}5}{0{,}08}\approx5{,}06.\]
		Vậy sau $6$ năm thì số tiền đầu tư ban đầu tăng thêm ít nhất  $50\%$.
	 }
\end{ex}

% nguồn: Ok_Toan11_TracNghiem_Lan2
\begin{ex}%[1D6H4-6]%[ma:43a4a4d1]
	Trong khuôn viên một trường đại học có $5\,000$ sinh viên, một  sinh viên vừa trở về sau kì nghỉ và bị nhiễm virus cúm truyền nhiễm kéo dài. Sự lây lan này được mô hình hóa bởi công thức
	$$y=\dfrac{5\,000}{1+4\,999\mathrm{e}^{-0{,}8t}},\;  t\ge0.$$
	Trong đó $y$ là tổng số sinh viên bị nhiễm sau $t$ ngày.  Các trường đại học sẽ cho các lớp nghỉ khi có nhiều hơn hoặc bằng $40\%$ số sinh viên bị lây nhiễm. Sau ít nhất bao nhiêu ngày thì trường cho các lớp nghỉ học?
	\choice 
		{\True $11$} 
		 {$12$}
		 {$10$}
		 {$13$}
	 \loigiai{
		Ta  có
		\allowdisplaybreaks 
		\begin{eqnarray*}
			&&\dfrac{5\,000}{1+4\,999\mathrm{e}^{-0{,}8t}}:5\,000\ge\dfrac{40}{100}\\ 
			 &\Leftrightarrow&1+4\,999\mathrm{e}^{-0{,}8t}\le\dfrac52\\
			&\Leftrightarrow&\mathrm{e}^{-0{,}8t}\le\dfrac{3}{9\,998}\\ 
			 &\Leftrightarrow&
			 t\ge\dfrac{\ln\left(\dfrac{3}{9\,998}\right)}{0{,}8}\approx10{,}14.
		\end{eqnarray*}
		Vậy sau  ít nhất $11$ ngày thì trường cho các lớp nghỉ học.
	} 
\end{ex}

% nguồn: Ok_TracNghiemLop11_Lan1
\begin{ex}%[1D6H4-6]%[ma:92a2f008]
	Sau khi bệnh nhân uống một liều thuốc, lượng thuốc còn lại trong cơ thể giảm dần và được tính theo công thức $P(t)=P_0 \cdot a^t(\mathrm{~mg})$, trong đó $a$ là hằng số dương, $t$ là thời gian tính bằng giờ kể từ thời điểm uống thuốc, $P_0(\mathrm{~mg})$ là lượng thuốc bệnh nhân đã uống. Biết rằng bệnh nhân đã uống $160 \mathrm{~mg}$ thuốc và sau $1$ giờ thì lượng  thuốc trong cơ thể còn lại $140 \mathrm{~mg}$. Sau $4$ giờ lượng thuốc đã giảm đi bao nhiêu phần trăm so với lượng thuốc ban đầu?
	\choice 
		{$93{,}27\%$}{$66{,}21\%$}{\True  $41{,}38\%$}{$58{,}62\%$}
	 \loigiai{
		Ta có $P(t)=P_0 \cdot a^t(\mathrm{~mg})$ nên $140=160\cdot  a^1\Leftrightarrow a=0{,}875$.\\
		$P(4)=160\cdot  {0{,}875}^4=\dfrac{12005}{128}$.\\
		Sau $4$ giờ số phần trăm lượng  thuốc đã giảm so với ban đầu.
		$$100-\left( \dfrac{12005}{128}:  160 \right)\cdot 100= 41{,}38\%.$$
		Do  đó Sau $4$ giờ lượng thuốc đã giảm đi $41{,}38\%$ so với lượng thuốc ban đầu.
	 }
\end{ex}

% nguồn: Ok_TracNghiemLop11_Lan1
\begin{ex}%[1D6H4-6]%[ma:dd698270]
	Biết rằng mức cường độ âm $L$ được tính bằng công thức $L=10\log\left(\dfrac{I}{I_0}\right)$ (dB), trong đó $I$ là cường độ của âm tính bằng W/m$^2$ và $I_0=10^{-12}$ W/m$^2$. Một giáo viên đang giảng bài trong lớp có mức cường độ âm bằng $55$ dB. Cường độ âm của giọng nói giáo viên  thuộc khoảng nào sau đây?
	 \choice
		 {\True $\left(10^{-7};10^{-6}\right)$}
		{$\left(10^{-8};10^{-7}\right)$} 
		 {$\left(10^{-5};10^{-4}\right)$}
		{$\left(10^{-6};10^{-5}\right)$} 
	 \loigiai{
		Ta có 
		\allowdisplaybreaks 
		\begin{eqnarray*}
			 L=10\log\left(\dfrac{I}{I_0}\right)&\Leftrightarrow& 55=10\log\left(\dfrac{I}{10^{-12}}\right)\\
			 &\Leftrightarrow&\log\left(\dfrac{I}{10^{-12}}\right)=5{,}5\\
			&\Leftrightarrow&\dfrac{I}{10^{-12}}=10^{5{,}5}\\ 
			 &\Leftrightarrow&I=10^{-6{,}5}\in \left(10^{-7};10^{-6}\right).
		\end{eqnarray*}
	 }
\end{ex}

% nguồn: Ok_TracNghiemLop11_Lan1
\begin{ex}%[1D6H4-6]%[ma:61587a45]
	Cho biết chu kì bán rã của chất phỏng xạ plutonium $Pu^{239}$  là $24.360$ năm (tức là lượng $Pu^{239}$ sau $24360$ năm phân hủy thì chỉ còn lại một nửa
	Sự phân hủy được tính bởi công thức $S=Ae^{rt}$, trong đó $A$ là lượng chất phóng xạ ban đầu, $r$ là tỉ lệ phân hủy hàng năm ($r < 0$), $t$ là thời gian phân huỷ, $S$ là lượng còn lại sau thời gian phân hủy $t$. Sau khoảng bao nhiêu năm  thì 10 gam $Pu^{239}$ phân hủy còn lại $1$ gam?
	 \choice
		{$833335$} 
		{\True  $822235$}
		 {$800235$}
		 {$822000$}
	 \loigiai{
		Vì chu kì bán rã của chất phỏng xạ  plutonium $Pu^{239}$ là $24.360$ năm nên ta có:
		\[5=10e^{r\cdot24360} \Leftrightarrow  r\approx-0,000028.
		\] 
		Sự phân hủy của chất phỏng xạ plutonium $Pu^{239}$ được tính bởi công  thức $S=Ae^{-0,000028t}$.\\
		Với $S=1$ thì $1=10e^{-0,000028t}$  $\Leftrightarrow t\approx 822235$ năm.\\
		Vậy sau khoảng $822235$ năm thì $10$  gam $Pu^{239}$ phân hủy còn lại $1$ gam.
	 }
\end{ex}

% nguồn: Ok_Toan11_TracNghiem_Lan2
\begin{ex}%[1D6V4-6]%[ma:047ed174]
	 Dân số thế giới được dự đoán theo công thức $P(t)=a\mathrm{e}^{bt}$, trong đó $a,b$ là các hằng số, $t$ là năm tính dân số. Theo số liệu thực tế, dân số thế giới năm $1950$ là $2\,560$ triệu người và năm $1980$ là $4\,440$ triệu người. Hãy dự đoán dân số thế giới năm $2030$ (làm tròn đến hàng triệu).
	 \choice
		{$19$ tỉ $280$  triệu}
		{$10$ tỉ $141$ triệu} 
		{$15$ tỉ $236$ triệu} 
		{\True $11$ tỉ $116$ triệu} 
	\loigiai{ 
		 Ta có
		 \[
		\heva{&a\mathrm{e}^{1\,950b}=2\,560,\\&a\mathrm{e}^{1\,980b}=4\,440} \Rightarrow \dfrac{a\mathrm{e}^{1\,980b}}{a\mathrm{e}^{1\,950b}}=\dfrac{4\,440}{2\,560} 
		 \Rightarrow
		 \mathrm{e}^{30b}=\dfrac{111}{64}.
		\] 
		 Khi đó
		\[ 
		 \dfrac{P(2\,030)}{P(1\,950)}
		=\dfrac{a\mathrm{e}^{2\,030}}{a\mathrm{e}^{1\,950b}} 
		 =\mathrm{e}^{80b}
		=\left(\mathrm{e}^{30b}\right)^{\tfrac83} 
		 =\left(\dfrac{111}{64}\right)^{\tfrac83}.
		 \]
		 Suy ra
		 \[
		P(21\,950\,030)=21\,950\,560\left(\dfrac{111}{64}\right)^{\tfrac83}\approx11\,116\ \text{(triệu người).} 
		 \]
	} 
\end{ex}


%%%%%%%%%%%%%%%% ĐÚNG/SAI %%%%%%%%%%%%%%%%

% nguồn: Ok_dungsai_lop11_Lan1
\begin{ex}%[1D6N4-6]%[ma:dd5fb00d]
 Các loài cây xanh trong quá trình quang hợp sẽ nhận được một lượng nhỏ cacbon $14$ (một đồng vị của cacbon). Khi một bộ phận của cây nào đó bị chết thì hiện tượng quang hợp cũng ngưng và nó sẽ không nhận thêm cacbon $14$ nữa. Lượng cacbon $14$ của bộ phận đó sẽ phân hủy một cách chậm chạp, nó chuyển hóa thành nitơ $14$. Gọi $P(t)$ là số phần trăm cacbon $14$ còn lại trong một bộ phận của cây sinh trưởng từ $t$ năm trước  đây thì $P(t)$ được tính theo công thức $P(t)=100\cdot (0,5)^{\frac{t}{5750}}$(\%). Phân tích một mẩu gỗ từ một công trình kiến trúc cổ, người ta thu được lượng cacbon $14$ còn lại trong mẫu gỗ đó là $50\%$. Gọi $a$ là niên đại của công trình kiến trúc.
  \choiceTF
   {\True $a>5000$}
   {\True $a<6000$}
  {\True $P(a)=50$} 
   {$P(a)=75$}
  \loigiai{
 Ta có $P(t)=100\cdot (0,5)^{\frac{t}{5750}}=50\Leftrightarrow t=5750$.\\ 
  Vậy niên đại công trình kiến trúc là $5750$ năm.
  }
\end{ex}

% nguồn: Ok_dungsai_lop11_Lan1
\begin{ex}%[1D6N4-6]%[ma:248ff7bf]
 Để biết dung dịch có tính axít, tính bazơ hay trung tính, người ta dùng độ $pH$ để xác định, biết $pH=-\log\left[H_3O^+\right]$. Trong đó, $pH$ là chữ đầu của nhóm từ ``potential of hydrogen'' nghĩa là tiềm lực của  hiđrô. Nếu $pH<7$: dung dịch có tính axít; $pH>7$: dung dịch có tính bazơ; $pH=7$: dung dịch trung tính. Hỏi nếu bia có nồng độ ion hiđrô $\left[H_3O^+\right]=0,00008$ thì
 \choiceTF 
  {\True Bia có tính axít} 
  {$pH>7$} 
  {\True  $pH \in (4;5)$}
  {$pH\in  (2;3)$}
  \loigiai{
 Từ giả thiết ta có độ  $pH$ của bia là $pH=-\log 0,00008\approx 4,1<7$.\\
  Do đó bia có tính axít.
 } 
\end{ex}

% nguồn: Ok_dungsai_lop11_Lan2
\begin{ex}%[1D6H4-6]%[ma:e2d0156d]
	Một người gửi số tiền $500$ triệu đồng vào ngân hàng với lãi suất $6{,}5$\% một năm  theo hình thức lãi kép.
	 \choiceTF
		{Lãi suất của ngân hàng là $0{,}65$  trong một năm}
		{\True Sau  khi gửi 1 năm, số tiền mà người đó có trong ngân hàng là $ 532\,500\,000$ đồng}
		{\True Sau khi gửi 3 năm, số tiền mà người đó có trong ngân hàng nhiều  hơn $ 600\,000\,000$ đồng}
		{Do thiếu tiền nên ở cuối năm thứ 3, người đó đã rút $100$ triệu đồng từ  ngân hang và tiếp tục gửi thêm 2 năm nữa thì rút toàn bộ số tiền. Lúc này người này có số tiền ít hơn $ 670\,000\,000$ đồng}
	\loigiai{ 
		\begin{itemchoice}
			\itemch Lãi  suất ngân hàng là $0{,}065$ trong một năm.
			\itemch  Sau một năm số tiền gửi là $ 500\left(1+6{,}5\%\right)^1=532{,}5$ (triệu đồng).
			\itemch Đến hết năm thứ ba, số tiền người đó có được là $ 500\left(1+6{,}5\%\right)^3>600$ triệu  đồng.
			\itemch Sau khi rút về 100 triệu đồng và tiếp tục gửi trong vòng 2 năm tiếp theo,  người đó có số tiền là $\left[500\left(1+6{,}5\%\right)^3-100\right]\left(1+6{,}5\%\right)^2\approx 571,621$ triệu đồng. \\
			Tổng số tiền người đó có được sau $5$ năm (sau khi làm tròn)  là $571,621+100=671,621$ triệu đồng, gần nhất với $671{,}620$ triệu đồng.
	\end{itemchoice}}
\end{ex}

% nguồn: Ok_dungsai_lop11_Lan2
\begin{ex}%[1D6H4-6]%[ma:e5df5b68]
	Lạm phát là sự tăng mức giá chung một cách liên tục của hàng hoá và dịch vụ theo thời gian, tức là sự mất giá trị của một loại tiền tệ  nào đó. Chẳng hạn, nếu lạm phát là $ 5$\% một năm thì sức mua của 1 triệu đồng sau một năm chỉ còn là 950 nghìn đồng (vì đã giảm mất $ 5$\% của 1 triệu đồng, tức là $50\,000$ đồng). Nói chung, nếu tỉ lệ lạm phát trung bình là $ r$\% một năm thì tổng số tiền $ P$ ban đầu, sau $ n$ năm số tiền đó chỉ còn giá trị là $ A=P\cdot{\left(1-\dfrac{r}{100}\right)^n}$.
	\choiceTF 
		{Nếu tỉ lệ lạm phát là  $6$\% một năm thì sức mua của $100$ triệu đồng sau một năm sẽ còn lại là $95$ triệu đồng}
		{\True Nếu tỉ lệ lạm phát là $8$\% một năm thì sức mua của $100$  triệu đồng sau hai năm sẽ còn lại là $84{,}64$ triệu đồng}
		{\True Nếu sức mua của $100$ triệu đồng sau hai năm chỉ còn là $90$ triệu đồng thì tỉ lệ lạm phát trung bình của hai năm đó là $5{,}13$\%} 
		{\True Nếu tỉ lệ lạm phát  là $5$\% một năm thì sau $14$ năm sức mua của số tiền ban đầu chỉ còn lại một nửa}
	 \loigiai{
		\begin{itemchoice}
			 \itemch Theo công thức $ A=P\left(1-\dfrac{r}{100}\right)^n$ ta có $ A=100\left(1-\dfrac{6}{100}\right)^1=94$ triệu đồng.\\
			Vậy sức mua của $100$ triệu đồng sau một năm  với tỉ lệ lạm phát là $ 6\%$ một năm chỉ còn lại khoảng $94$ triệu đồng.
			\itemch Theo công thức $ A=P.\left(1-\dfrac{r}{100}\right)^n$ ta có $ A=100.\left(1-\dfrac{8}{100}\right)^2=84,64$  triệu đồng\\
			Vậy sức mua của $100$ triệu đồng sau hai năm với tỉ  lệ lạm phát là $ 8\%$ một năm chỉ còn lại khoảng 84,64 triệu đồng.
			\itemch Thay $ P=100$ triệu đồng, $ A=90$ triệu đồng, $ n=2$ vào phương trình ta có 
			$$90=100\left(1-\dfrac{r}{100}\right)^2\Rightarrow r=5{,}13\%.$$ 
			Vậy  tỉ lệ lạm phát trung bình của hai năm đó là khoảng $5{,}13$\%.
			\itemch Thay $ P=1$ và  $ A=\dfrac{1}{2}$ vào phương trình ta có $\dfrac{1}{2}=\left(1-\dfrac{r}{100}\right)^n$.
			 $$\ln\left(\dfrac{1}{2}\right)=n\ln\left(1-\dfrac{r}{100}\right)\Leftrightarrow n=\dfrac{\ln\left(\dfrac{1}{2}\right)}{\ln\left(1-\dfrac{r}{100}\right)}\Leftrightarrow n=\dfrac{\ln\left(\dfrac{1}{2}\right)}{\ln\left(1-\dfrac{5}{100}\right)}\approx 13{,}51.$$
			Vậy sau khoảng 14 năm thì sức mua của số tiền ban đầu sẽ chỉ còn lại một nửa nếu tỉ lệ lạm phát là $ 5\%$ một năm. 
		\end{itemchoice}
	} 
\end{ex}

% nguồn: Ok_dungsai_lop11_Lan2
\begin{ex}%[1D6H4-6]%[ma:31a67e2d]
	 Anh Nam vay tiền ngân hàng $1$ tỷ đồng theo phương thức trả góp với lãi suất $0{,}5$\%/ tháng. Nếu cuối mỗi tháng bắt đầu từ tháng thứ nhất anh Nam trả $ 30$ triệu đồng. Biết rằng lãi suất ngân hàng không thay đổi trong suốt thời gian trả nợ.
	 \choiceTF
		{\True Số tiền nợ sau $ 8$  tháng là $796\,464\,780{,}4$}
		 {\True Số tiền nợ sau $ 10$ tháng là $744\,299\,339,8$}
		{\True Sau  $ 37$ tháng thì anh Nam trả hết nợ}
		{\True Nếu cuối mỗi tháng  bắt đầu từ tháng thứ nhất anh Nam trả $ 45$ triệu đồng thì sau hai năm anh Nam trả hết nợ}
	 \loigiai{
		\begin{itemchoice}
			\itemch Gọi $a$ là số tiền vay, $r$ là lãi  suất, $ m$ là số tiền hàng tháng trả.\\
			 Số tiền nợ sau tháng thứ nhất là $N_1=a\left(1+r\right)-m$.\\
			Số tiền nợ sau tháng thứ hai là  $$N_2=\left[a\left(1+r\right)-m\right]+\left[a\left(1-r\right)-m\right]r-m=a{\left(1+r\right)^2}-m\left[\left(1+r\right)+1\right]$$\\
			 ….\\
			Số  tiền nợ sau $ n$ tháng là $$N_n=a{\left(1+r\right)^n}-m\left[\left(1+r\right)^{n-1}+\left(1+r\right)^{n-2}+...+1\right]=a{\left(1+r\right)^n}-m\dfrac{\left(1+r\right)^n-1}{r}.$$
			Sau $ n$ tháng  anh Nam còn nợ số tiền là $$N_n=10^9\left(1+0{,}5\%\right)^n-30\cdot10^6\dfrac{\left(1+0{,}5\%\right)^n-1}{0{,}5\%}.$$
			Số tiền nợ  sau $ 8$ tháng là $796\,464\,780{,}4$.
			\itemch Số  tiền nợ sau $ 10$ tháng là $744\,299\,339,8$.
			\itemch Sau $ n$ tháng anh Nam trả  hết nợ: $N_n=a{\left(1+r\right)^n}-m\dfrac{\left(1+r\right)^n-1}{r}=0$$\Leftrightarrow 1000\left(1+0{,}005\right)^n-30\dfrac{\left(1+0{,}005\right)^n-1}{0{,}005}=0\Leftrightarrow n=36{,}55$\\
			Vậy $ 37$ tháng  thì anh Nam trả hết nợ.
			\itemch  Sau $ n$ tháng anh Nam trả hết nợ $$N_n=a{\left(1+r\right)^n}-m\dfrac{\left(1+r\right)^n-1}{r}=0\Leftrightarrow 1000\left(1+0{,}005\right)^n-45\cdot\dfrac{\left(1+0{,}005\right)^n-1}{0{,}005}=0\Leftrightarrow n=23{,}615.$$
			Vậy sau  hai năm anh Nam trả hết nợ.
	\end{itemchoice}}
\end{ex}

% nguồn: Ok_dungsai_lop11_Lan2
\begin{ex}%[1D6H4-6]%[ma:cef238bb]
	Trong năm $2\,019$, diện tích rừng trồng mới của tỉnh A là $1\,000$ ha. Giả sử diện tích rừng trồng mới của tỉnh A mỗi năm tiếp theo đều tăng $6$\% so với diện tích rừng trồng mới  của năm liền trước. Sau năm $2\,019$.
	\choiceTF 
		{Công thức sau $n$ năm  thì diện tích trồng rừng của tỉnh A là $ A=1000\left(1+0{,}06\right)^{n+1}$}
		{\True Vào năm $2032$, diện tích rừng năm đó hơn gấp  đôi năm $2019$}
		{\True Vào năm $2025$ thì diện tích rừng  năm đó đạt trên $1\,400$ ha}
		{\True  Diện tích rừng vào hai năm sau kể từ năm $2019$ sẽ đạt $1\,123{,}6$ ha}
	\loigiai{ 
		Áp dụng  công thức $A=a{\left(1+r\right)^n}$.
		\begin{itemchoice}
			\itemch Ta có sau $ n$ năm thì diện  tích trồng mới của tỉnh A là $ A=1000\left(1+0{,}06\right)^n$.
			\itemch $ 1000\cdot\left(1+0{,}06\right)^n>2000\Rightarrow  1,06^n>1\Rightarrow n>11,8$.\\
			Vào năm $2032$, diện tích rừng năm đó hơn  gấp đôi năm $2019$.
			\itemch Theo đề bài,  ta có $ 1000\cdot\left(1+0{,}06\right)^n>1400\Rightarrow 1{,}06^n>1{,}4\Rightarrow n>5{,}774$\\
			Vậy vào  năm 2025 thì diện tích rừng mới năm đó đạt trên $1\,400$ ha.
			\itemch $n=2\Rightarrow 1000\left(1+0{,}06\right)^2=1\,123{,}6$ ha. 
	\end{itemchoice}}
\end{ex}

% nguồn: Ok_dungsai_lop11_Lan1
\begin{ex}%[1D6V4-6]%[ma:37cfde0c]
 Một bà mẹ Việt Nam anh hùng được hưởng số tiền là $4$ triệu đồng mỗi tháng (chuyển vào tài khoản ngân hàng của mẹ ở ngân hàng vào đầu mỗi tháng).  Tháng $12$ năm $2023$ về trước, mẹ đã rút hết tiền hàng tháng. Từ tháng $1$ năm $2024$, mẹ không đi rút tiền mà để lại ngân hàng và được tính lãi $1$\%/tháng vào ngày mồng $1$ đầu tháng trên tổng số tiền gốc và lãi có của tháng liền trước đó.
 \choiceTF 
  {Đến  ngày $29$ tháng $2$ năm $2024$, mẹ đi rút tiền thì tổng số tiền nhận được bằng $8\,080\,000$ đồng}
  {\True Đến ngày $02$ tháng $3$ năm $2024$, mẹ đi rút tiền gốc và lãi của tháng $1$, tháng $2$ và tiền tháng $3$, thì tổng số tiền nhận được bằng  $12\,120\,400$ đồng}
  {\True Đến ngày $05$  tháng $01$ năm $2025$, mẹ đi rút tiền thì số tiền nhận được bằng $50$ triệu $730$ nghìn đồng (làm tròn đến hàng nghìn)}
  {Để nhận được số tiền hơn $100$ triệu đồng thì mẹ phải đi  rút tiền trong tháng $12$ năm $2025$}
 \loigiai{ 
  \begin{itemchoice}
   \itemch  \textbf{Sai.}\\
    Đến ngày $29$ tháng $2$ năm $2024$, mẹ đi rút tiền thì số tiền nhận được bao gồm:
   Tiền tháng $1$ và Lãi, cộng với tiền tháng $2$ chưa  có lãi, nên số tiền nhận được bằng:
    \[
   4\,000\,000\cdot(1+1\%)+4\,000\,000=8\,040\,000\ \text{đồng}. 
   \] 
   \itemch  \textbf{Đúng.}\\
   Đến ngày $02$ tháng $3$ năm $2024$, mẹ đi rút tiền thì số tiền nhận được bao gồm tiền tháng $1$, tháng $2$ và lãi, cộng với tiền tháng $3$ chưa có lãi, nên  số tiền nhận được bằng
    \[(4\,000\,000\cdot 1{,}01+4\,000\,000) \cdot 1{,}01+4\,000\,000=12\,120\,400 \ \text{đồng}.\]
   \itemch  \textbf{Đúng.}\\
   Xét bài toán  tổng quát:
   Gọi số tiền mẹ nhận  được gửi vào ngân hàng mỗi đầu hàng tháng là $A$ đồng.
   Lãi suất hàng tháng mẹ gửi tại ngân  hàng là $r \%$.\\
    Đầu tháng thứ $2$, số tiền của mẹ là $A+A(1+r)$ đồng.\\
   Đầu  tháng thứ $3$, số tiền của mẹ là
   \[ 
   A+[A+A\cdot  (r+1)](1+r)=A+A(1+r)+A(1+r)^2\ \text{đồng}.
   \] 
   Đầu  tháng thứ $n$, số tiền của mẹ là
   \[A+A(1+r)+A(1+r)^2+\cdots+A(1+r)^{n-1}\  \text{đồng}.\]
   Ta có 
   \[A+A(1+r)+A(1+r)^2+\cdots+A(1+r)^{n-1}=A\cdot \dfrac{(1+r)^n-1}{(1+r)-1}=A\cdot \dfrac{(1+r)^n-1}{r}.\] 
   Đến ngày $05$ tháng $1$ năm $2025$ thì $n=12$, với $A=4\,000\,000$; $r=1\%=0{,}01$, số tiền nhận được  bằng
   $4\,000\,000\cdot \dfrac{1{,}01^{12}-1}{0{,}01}  \approx 50\,730\,012$ đồng, làm tròn là $50$ triệu $730$ nghìn đồng.
    \itemch \textbf{Sai.}\\
   Để nhận được số tiền  hơn $100$ triệu đồng thì
    \allowdisplaybreaks
   \begin{eqnarray*}
     && 4\,000\,000\cdot \dfrac{1{,}01^n-1}{0{,}01} \geq 100\,000\,000\\
    &\Rightarrow& 1{,}01^n \geq \dfrac{5}{4}  \\
    &\Rightarrow& n \geq \log_{1{,}01} \dfrac{5}{4} \approx 22{,}4. 
   \end{eqnarray*}
   Vậy phải $23$ tháng, tức là đến tháng $11$ năm $2025$ thì nhận được  số tiền hơn $100$ triệu đồng.
  \end{itemchoice}
  }
\end{ex}

% nguồn: Ok_dungsai_lop11_Lan1
\begin{ex}%[1D6C4-6]%[ma:c2063d75]
 Lạm phát là sự tăng mức giá chung một cách liên  tục của hàng hoá và dịch vụ theo thời gian, tức là sự mất giá trị của một loại tiền tệ nào đó. Chẳng hạn, nếu lạm phát là $5 \%$ một năm thì sức mua của $1$ triệu đồng sau một năm chỉ còn là $950$ nghìn đồng. Nói chung, nếu tỉ lệ lạm phát trung bình là $r \%$ một năm thì tổng số tiền $P$ ban đầu, sau $n$ năm số tiền đó chỉ còn giá trị là $A=P\left(1-r \% \right)^n$. Các mệnh đề sau Đúng hay Sai?
 \choiceTF 
  {\True Nếu tỉ lệ lạm phát  là $7 \%$ một năm thì sức mua của $100$ triệu đồng sau hai năm sẽ còn lại là $86\,490\,000$ đồng}
  {Nếu tỉ lệ lạm phát là $7 \%$ một năm thì sức mua của $100$ triệu đồng sau hai năm sẽ còn lại  $96\,490\,000$ đồng}
  {Nếu sức mua của $100$ triệu đồng sau  ba năm chỉ còn lại $80$ triệu đồng thì tỉ lệ lạm phát trung bình của ba năm đó là $9{,}17 \%$}
  {Nếu tỉ lệ lạm phát trung bình là $6 \%$ một năm thì sau ít nhất $15$ năm sức mua của số tiền ban đầu  chỉ còn lại không quá một nửa}
 \loigiai{ 
  \begin{itemchoice}
   \itemch Đúng. Sau $2$ năm, sức mua còn lại của $100$ triệu đồng  là \[A=10^8\cdot\left(1-7 \% \right)^2=86\,490\,000\text{ (đồng).}\]
   \itemch  Sai.
   \itemch Sai. Ta có $10^8\cdot\left(1-r \%  \right)^3=8\cdot 10^7 \Rightarrow r\approx 7{,}17$.
   \itemch Sai. Gọi số tiền ban đầu là $P$,  ta có $P\cdot\left(1-6 \% \right)^n=\dfrac{P}{2} \Rightarrow n\approx 11{,}2$.\\
   Vậy sau ít nhất $12$ năm thì sức mua của số tiên ban đầu chỉ còn lại  không quá một nửa.
  \end{itemchoice}
  }
\end{ex}

% nguồn: Ok_dungsai_lop11_Lan1
\begin{ex}%[1D6C4-6]%[ma:1a3ca92c]
 Ông Xuân gửi $200$ triệu đồng vào ngân hàng theo phương thức lãi kép với  lãi suất cố định là $0,5\%$ một tháng. Cùng thời điểm đó, ông Thắng gửi $10$ triệu đồng mỗi tháng vào ngân hàng đó theo phương thức lãi kép với lãi suất cố định là $0{,}5\%$ một tháng. (làm tròn kết quả quả đến hàng nghìn)
 \choiceTF 
  {\True Số tiền cả gốc và lãi ông Xuân nhận được sau $10$ tháng  là $210\,228\,000$}
   {Sau ít nhất $81$ tháng thì số tiền của ông Xuân nhiều hơn $300$ triệu đồng}
  {\True Sau $1$ năm số tiền cả  gốc và lãi trong ngân hàng của ông Thắng là $123\,972\,000$}
  {\True Sau ít nhất $22$ tháng thì số tiền của ông Thắng lớn hơn số  tiền của ông Xuân}
 \loigiai{ 
  \begin{itemchoice}
   \itemch Đúng. Số tiền cả gốc và lãi ông Xuân nhận được sau $10$ tháng là  \[200\,000\,000\left(1+\dfrac{1}{200} \right)^{10} \approx 210\,228\,026.\]
   \itemch Sai. Theo  công thức lãi kép, số tiền cả gốc và lãi ông Xuân nhận được sau $n$ tháng là \[200\,000\,000\left(1+\dfrac{1}{200} \right)^n\]
   Do đó $200\,000\,000\left(1+\dfrac{1}{200} \right)^n > 300\,000\,000\Leftrightarrow n > \log _{\frac{201}{200}} \dfrac{3}{2} \approx 81{,}29558565$. 
   \itemch  Đúng. Theo công thức lãi kép,
   \begin{itemize}
    \item Số tiền cả gốc và  lãi ông Thắng nhận được sau $1$ tháng là
    \[10\,000\,000\left(1+\dfrac{1}{200} \right)^1. 
    \] 
    \item Số tiền cả gốc và lãi  ông Thắng nhận được sau $2$ tháng là
    \begin{eqnarray*}
     & &10\,000\,000\left(1+\dfrac{1}{200} \right)^2+10\,000\,000\left(1+\dfrac{1}{200}  \right)^1\\
     &=&10\,000\,000\left[\left(1+\dfrac{1}{200} \right)^1+\left(1+\dfrac{1}{200} \right)^2\right]. 
    \end{eqnarray*}
    \item Số tiền cả gốc và lãi ông Thắng nhận được  sau $3$ tháng là
    \begin{eqnarray*}
      & &10\,000\,000\left(1+\dfrac{1}{200} \right)^3+10\,000\,000\left(1+\dfrac{1}{200} \right)^2+10\,000\,000\left(1+\dfrac{1}{200} \right)^1\\
     &=&10\,000\,000\left[\left(1+\dfrac{1}{200} \right)^1+\left(1+\dfrac{1}{200}  \right)^2+\left(1+\dfrac{1}{200} \right)^3\right].
    \end{eqnarray*}
     \dotfill
     \item Số tiền sau $1$ năm là
    \begin{eqnarray*}
     & &10\,000\,000\left[\left(1+\dfrac{1}{200}  \right)^1+\left(1+\dfrac{1}{200} \right)^2+\cdots+\left(1+\dfrac{1}{200} \right)^{12} \right]\\
     &=&10\,000\,000\dfrac{\left(1+\dfrac{1}{200} \right)\left[\left(1+\dfrac{1}{200}  \right)^{12}-1\right]}{\dfrac{1}{200}}\\
      &\approx& 123\,972\,401.
    \end{eqnarray*}
   \end{itemize}
   \itemch Đúng. Theo  công thức lãi kép, số tiền cả gốc và lãi ông Xuân và ông Thắng nhận được sau $n$ tháng lần lượt là
   \begin{eqnarray*}
    & &200\,000\,000\left(1+\dfrac{1}{200} \right)^n;\; 10\,000\,000\dfrac{\left(1+\dfrac{1}{200} \right)\left[\left(1+\dfrac{1}{200} \right)^n-1\right]}{\dfrac{1}{200}}\\ 
    &=&2\,010\,000\,000\left[\left(1+\dfrac{1}{200}  \right)^n-1\right].
   \end{eqnarray*}
   \[ 
   \] 
   Do  đó
   \begin{eqnarray*}
     & &2\,010\,000\,000\left[\left(1+\dfrac{1}{200} \right)^n-1\right] > 200\,000\,000\left(1+\dfrac{1}{200} \right)^n\\
    &\Leftrightarrow& n > \log _{\frac{201}{200}} \dfrac{201}{181} \approx 21{,}01393574. 
   \end{eqnarray*}
  \end{itemchoice}
  }
\end{ex}

% nguồn: Ok_dungsai_lop11_Lan1
\begin{ex}%[1D6C4-6]%[ma:9e83741f]
 Ông A đem $800$ triệu đồng gửi vào một ngân hàng với lãi suất $0{,}5\%$ một tháng. Biết rằng nếu không rút tiền ra khỏi ngân hàng thì cứ sau mỗi tháng số tiền lãi sẽ được nhập vào gốc để tính lãi cho tháng tiếp theo và từ tháng thứ hai trở đi. Lãi suất được cho là không đổi trong suốt  thời gian vậy tiền.
 \choiceTF 
   {Số tiền cả gốc và lãi ông A rút về sau một năm lớn hơn $850$ triệu đồng}
  {\True Ông A định dùng tiền lãi sau $2$ năm để mua chiếc xe SH trị giá $100$ triệu đồng. Sau Đúng  $2$ năm tiền lãi thu được đủ để ông A mua chiếc xe đó}
  {\True Sau ít nhất $45$ tháng thì số tiền thu về cả gốc lẫn lãi lớn  hơn $1$ tỷ đồng}
  {\True Sau khi gửi, cứ tròn mỗi tháng, ông đến ngân hàng  rút $6$ triệu để chi tiêu cho gia đình. Một năm sau khi rút tiền, số tiền tiết kiệm của ông An còn lại ít hơn $776$ triệu đồng}
  \loigiai{
  \begin{itemchoice}
   \itemch Sai. Số tiền cả  gốc và lãi ông A rút về sau một năm
   \[A_{12}=A\left(1+r\right)^{12}=800\,000\,000\cdot\left(1+0{,}5\%  \right)^{12} \approx 849\,342\,250 < 850\,000\,000.
    \]
   \itemch  Đúng. Số tiền lãi ông A rút về sau hai năm \[A_{12}=A\left(1+r\right)^{12}=800\,000\,000\cdot\left(1+0{,}5\% \right)^{24}-800\,000\,000\approx 101\,728\,000> 100\,000\,000.\]
   \itemch Đúng. Số tiền  thu về cả gốc lẫn lãi lớn hơn $1$ tỷ đồng
   \begin{eqnarray*}
    800\,000\,000\cdot \left(1+0{,}5\% \right)^n &>& 1\,000\,000\,000\\ \Leftrightarrow \left(1+0{,}5\% \right)^n &>& \dfrac{5}{4}\\ \Leftrightarrow n  &>& \log _{\left(1+0{,}5\% \right)} \left(\dfrac{5}{4} \right)\\
    \Leftrightarrow n &>&  44{,}74\text{ tháng}.
   \end{eqnarray*}
   \itemch Đúng. Gửi ngân hàng số tiền là A đồng với lãi suất $r$ /tháng.\\ Mỗi tháng vào ngày ngân hàng tính lãi, rút ra số  tiền là $X$ đồng.\\
   Sô tiền còn lại sau $n$ tháng  đươc tính theo công thức
   \begin{eqnarray*}
    S_n&=&A\left(1+r\right)^n-X\dfrac{\left(1+r\right)^n-1}{r}\\ 
     &=&800\left(1{,}005\right)^{12}-6\cdot\dfrac{\left(1{,}005\right)^{12}-1}{0{,}005}\\
     &=&1200-400\cdot(1{,}005)^{12}\\
     &=&775{,}3288753.
   \end{eqnarray*}
  \end{itemchoice}
 } 
\end{ex}

% nguồn: Ok_dungsai_lop11_Lan1
\begin{ex}%[1D6C4-6]%[ma:dde6bd19]
 Một người gửi $100$ triệu đồng  vào một ngân hàng với lãi suất $0{,}6\%$ /tháng. Biết rằng nếu không rút tiền ra khỏi ngân hàng thì cứ sau mỗi tháng, số tiền lãi sẽ được nhập vào vốn để tính lãi cho tháng tiếp theo và trong suốt thời gian gửi tiền người đó không rút tiền và lãi suất không thay đổi.
 \choiceTF 
  {Sau  $1$ tháng, số tiền cả vốn lẫn lãi là $106$ triệu đồng}
  {Sau $2$ tháng, số tiền lãi  thu được là $1\,203\,600$ đồng}
  {Sau $n$ tháng, số tiền cả vốn lẫn lãi là $100\cdot\left(1+0{,}6\% \right)^{n+1}$} 
  {\True Để thu được nhiều hơn $10$ triệu tiền lãi, cần phải gửi ít  nhất $16$ tháng}
  \loigiai{
  \begin{itemchoice}
   \itemch Sau $1$ tháng, số  tiền cả vốn lẫn lãi là $100\cdot\left(1+0{,}6\% \right)^1=100{,}6$ (triệu đồng).
   \itemch Sau $1$ tháng, số tiền  cả vốn lẫn lãi là $100\cdot\left(1+0{,}6\% \right)^2=101{,}2036$ (triệu đồng).
   \itemch Sau $n$ tháng, người đó lĩnh được số tiền  là $100\cdot\left(1+0{,}6\% \right)^n$.
   \itemch Ta có $100\cdot\left(1+0{,}6\% \right)^n-100> 10\Leftrightarrow n> \log _{1+0{,}6\%}  \dfrac{11}{10} \approx 15{,}9$ (triệu đồng).\\
   Vậy để thu được nhiều hơn $10$ triệu tiền lãi, cần phải gửi ít nhất  $16$ tháng
  \end{itemchoice}
  }
\end{ex}

% nguồn: Ok_dungsai_lop11_Lan1
\begin{ex}%[1D6C4-6]%[ma:d9afeee5]
 Cent âm nhạc là một đơn vị trong thang lôgarit của cao độ hoặc khoảng tương đối. Một quãng  tám bằng $1\,200$ cent. Công thức xác định chênh lệch khoảng thời gian giữa hai nốt nhạc có tần số $a$ và $b$ ($a > b$) là $n=1200\cdot\log_2 \dfrac{a}{b}$.
  \choiceTF
  {Khoảng thời gian chênh lệch giữa hai nốt nhạc có tần số $443$ Hz và $415$ Hz là  $131$ cent}
  {\True Khoảng thời gian chênh lệch giữa hai nốt  nhạc có tần số $345$ Hz và $398$ Hz nằm trong khoảng $\left(246;250\right)$}
   {\True Giả sử khoảng thời gian là $230$ cent và tần số đầu là $328$ Hz thì tần số cuối cùng là $287{,}2$ Hz}
  {Với tần số đầu không vượt quá $355$ Hz  và tần số cuối cùng là $384$ Hz thì khoảng thời gian chênh lệnh giữa hai nốt nhạc không vượt quá $178$ cent}
  \loigiai{
  \begin{itemchoice}
   \itemch Sai. Khoảng thời gian chênh lệch giữa hai nốt nhạc có tần số $443$ Hz  và $415$ Hz là \[n=1\,200\cdot\log_2 \dfrac{443}{415} \approx 113\,\mathrm{cent}.\]
   \itemch Đúng. Khoảng thời  gian chênh lệch giữa hai nốt nhạc có tần số $345$ Hz và $398$ Hz là
   \[n=1200\cdot\log_2  \dfrac{398}{345} \approx 247{,}4.
    \]
   \itemch Đúng. Ta có phương trình \[230=1\,200\cdot\log_2 \dfrac{328}{b}  \Leftrightarrow \log_2 \dfrac{328}{b}=\dfrac{23}{120} \Leftrightarrow \dfrac{328}{b}=2^{\tfrac{23}{120}} \Leftrightarrow b\approx 287{,}2\,\mathrm{Hz}.\]
   \itemch Sai. Ta có phương trình \[178\ge 1\,200\cdot\log_2 \dfrac{a}{320} \Leftrightarrow \log_2 \dfrac{a}{320} \le \dfrac{89}{600} \Leftrightarrow \dfrac{a}{320} \le 2^{\frac{89}{600}} \Leftrightarrow  a\le 354{,}6\,\mathrm{Hz}.\]
  \end{itemchoice}
 } 
\end{ex}

% nguồn: Ok_dungsai_lop11_Lan1
\begin{ex}%[1D6C4-6]%[ma:e3064423]
 Một người gửi ngân hàng  $100$ triệu theo thể thức lãi kép, lãi suất $0{,}5\%$ một tháng. Mỗi mệnh đề sau đây Đúng hay Sai?
 \choiceTF 
  {\True Sau $3$ tháng số tiền cả gốc và lãi người đó nhận được nhiều hơn $101$ triệu đồng} 
  {Sau ít nhất $40$ tháng, người đó nhận được tổng số tiền nhiều hơn $125$ triệu} 
  {Số tiền lãi thu được sau $3$ tháng khi gửi lãi suất $0,5\%$ một tháng nhiều hơn số tiền lãi  thu được sau $4$ tháng nếu gửi lãi suất $0{,}4\%$ một tháng}
  {Tổng số tiền gốc và lãi gấp đôi số  tiền ban đầu sau $10$ năm}
 \loigiai{ 
  \begin{itemchoice}
   \itemch Đúng. Sau  $3$ tháng số tiền cả gốc và lãi là $P_3=100(1+0{,}005)^3 \approx 101\,507\,513$.
   \itemch Sai. Sau $n$ tháng, người đó nhận được $P_n=100(1+0{,}005)^n$ triệu đồng.\\ 
   Theo đề bài \[P_n=100(1+0{,}005)^n > 125\Leftrightarrow (1{,}005)^n > 1{,}25\Leftrightarrow n  > \log _{1\,005} 1{,}25=44{,}7\text{ (tháng).}\]
   Vậy sau $45$ tháng, người đó có nhiều hơn $125$  triệu đồng.
   \itemch  Sai. Sau $3$ tháng số tiền cả gốc và lãi là $P_3=100(1+0{,}005)^3 \approx 101\,507\,513$.\\
   Số tiền lãi là $1\,507\,513$ đồng.\\ 
   Số tiền gốc và lãi thu được sau $4$ tháng nếu gửi lãi suất $0{,}4\%$  một tháng
   \[A_4=100(1+0{,}004)^4 \approx  101\,609\,625.
   \] 
   Số tiền lãi là  $1\,609\,625$.
   \itemch Sai. $200=100(1+0{,}005)^n \Leftrightarrow (1{,}005)^n=2\Leftrightarrow n=\log _{1\,005}  2=138{,}975$.\\
   Sau số năm là $138{,}975:12=11{,}58>10$ năm. 
  \end{itemchoice}
 } 
\end{ex}

% nguồn: Ok_dungsai_lop11_Lan1
\begin{ex}%[1D6C4-6]%[ma:146b81d4]
 Người ta dùng thuốc để khử khuẩn cho một thùng nước. Biết rằng nếu lúc đầu mỗi mililít nước chứa $P_0$ vi khuẩn thì sau $t$ giờ, số lượng vi khuấn trong mỗi mililít nước là $P=P_0 \cdot 10^{-\alpha t}$, với $\alpha$ là một hằng số dương nào đó. Biết rằng ban đầu mỗi mililit nước có $4\,000$  vi khuẩn và sau $2$ giờ, số lượng vi khuẩn trong mỗi mililít nước là $1\,000$. Xét tính Đúng Sai các mệnh đề sau.
  \choiceTF
  {$\alpha$ nằm trong  khoảng $\left(1;2\right)$}
  {\True Sau $3$ giờ $30$ phút thì lượng vi khuẩn trong mỗi  mililít nước ít hơn $500$}
  {\True Lượng vi khuẩn mất đi trong mỗi mililít trong khoảng thời gian từ $1$ giờ đến $2{,}5$ giờ  tính từ lúc dùng thuốc thì lớn hơn $1\,200$}
  {\True Lượng vi khuẩn sau khoảng $1{,}32$ giờ sẽ bằng $40\%$ lượng vi khuẩn  ban đầu}
  \loigiai{
  \begin{itemchoice}
   \itemch Sai.  Ta có
   \begin{eqnarray*}
     P=P_0 \cdot 10^{-\alpha t}&\Rightarrow& 1\,000=4\,000\cdot 10^{-2\alpha}\\
    &\Rightarrow& \alpha=-\dfrac{1}{2} \log \dfrac{1}{4}=\dfrac{1}{2} \log 4\approx  0{,}31\notin \left(1;2\right).
   \end{eqnarray*}
   \itemch  Đúng.
   \begin{eqnarray*}
    P=P_0\cdot 10^{-\alpha t} \le 500&\Leftrightarrow& 4\,000\cdot 10^{-\alpha  t} \le 500\\
    &\Leftrightarrow& 10^{-\alpha t} \le \dfrac{1}{8}\\ 
    &\Leftrightarrow&-\alpha  t\le \log \dfrac{1}{8}\\
    &\Leftrightarrow& t\ge-\dfrac{1}{\alpha} \log \dfrac{1}{8}\\ &\Leftrightarrow& t\ge \dfrac{-1}{\dfrac{1}{2} \log 4} \cdot \log \dfrac{1}{8}\\ 
     &\Leftrightarrow& t\ge 3.
   \end{eqnarray*}
   \itemch Đúng. Lượng vi khuẩn trong mỗi mililít  nước sau $1$ giờ là $4\,000\cdot 10^{-\frac{1}{2} \log 4\cdot 1}=2\,000$.\\
   Lượng vi khuẩn trong mỗi mililít nước sau $2{,}5$ giờ là $4\,000\cdot  10^{-\frac{1}{2} \log 4\cdot 2{,}5} \approx 707$.\\
   Lượng vi khuẩn mất đi khoảng  $2\,000-707=1\,293> 1\,200$.
   \itemch Đúng. $P=P_0 \cdot  10^{-\alpha t}=\dfrac{40}{100}\cdot P_0 \Leftrightarrow 10^{-\alpha t}=\dfrac{2}{5} \Leftrightarrow t=\dfrac{\log \tfrac{2}{5}}{-\alpha}=\dfrac{\log \tfrac{2}{5}}{-\tfrac{1}{2} \log 4} \approx 1{,}32$.
  \end{itemchoice}
  }
\end{ex}

% nguồn: Ok_dungsai_lop11_Lan1
\begin{ex}%[1D6C4-6]%[ma:c98a5465]
 Để đặc trưng cho độ to nhỏ của âm, người ta đưa ra khái niệm mức cường độ của âm. Một đơn vị thường dùng để đo mức cường độ của âm là đề-xi-ben. Khi đó mức cường độ  $L$ của âm được tính theo công thức $L=10\log \dfrac{I}{I_0}$. Trong đó, $I$ là cường độ của âm tại thời điểm đang xét, $I_0$ cường độ âm ở ngưỡng nghe ($I_0=10^{-12}$ (w/m$^2$)).
 \choiceTF 
  {Một cuộc trò chuyện bình thường  trong lớp học có mức cường độ âm trung bình là $68$ dB. Cường độ âm tương ứng ra lớn hơn $6{,}5\cdot 10^{-6}$ (w/m$^2$)}
  {\True Hai cây đàn ghi ta giống nhau, cùng hòa tấu một bản nhạc.  Mỗi chiếc phát ra âm có mức âm trung bình là $60$ dB. Mức cường độ âm tổng cộng do hai chiếc đàn cùng phát ra xấp xỉ $63$ dB}
  {Tiếng ồn phát ra ở xưởng cưa ở mức cường độ âm đo được là $93$ dB do $7$ chiếc cưa máy giống nhau  cùng hoạt động gây ra. Giả sử có $3$ chiếc cưa máy đột ngột ngừng hoạt động thì mức cường độ âm trong xưởng lúc này nhỏ hơn $90$ dB. }
  {Tiếng ồn phát ra từ tiếng phím liên tục ở một bàn phím  của máy tính có cường độ âm đo được là $10^{-5}$ (w/m$^2$). Giả sử phòng làm việc của một công ty có hai nhân viên văn phòng cùng thực hiện thao tác gõ phím trên hai bàn phím máy tính giống nhau thì mức cường độ âm tổng cộng đo cả hai bàn phím phát ra cùng lúc là $70$ dB}
 \loigiai{ 
  \begin{itemchoice}
   \itemch Sai.  vì theo giả thiết ta có $L=68$ dB, $I_0=10^{-12}$ (w/m$^2$). Tính $I$.\\
   Áp dụng công  thức ta có
   \begin{eqnarray*}
    & &L=10\log \dfrac{I}{I_0} \Leftrightarrow 68=10\log \dfrac{I}{I_0}\\ 
    &\Leftrightarrow& \dfrac{I}{I_0}=10^{6{,}8} \Leftrightarrow I=10^{6{,}8}\cdot 10^{-12} \approx 6{,}3\cdot10^{-6}\mathrm{w/m^2}. 
   \end{eqnarray*}
   \itemch Đúng. Vì Mức cường độ âm do một chiếc đàn ghita phát ra là $L=10\log  \dfrac{I}{I_0}=60$.\\
   Mức cường độ âm  do hai chiếc đàn ghita phát ra là \[L_2=10\log \dfrac{2I}{I_0}=10\cdot\left(\log 2+\log \dfrac{I}{I_0} \right)=10\log 2+10\log \dfrac{I}{I_0}=10\log 2+60\approx 63\left(\mathrm{dB}\right).\]
   \itemch Sai. Vì gọi cường độ của âm do  $1$ cái cưa phát ra là $I_1$.\\
   Lúc đầu mức cường độ âm là  $7I_1$.
    \[L=10\log \dfrac{7I_1}{I_0}=93\Leftrightarrow 10\log 7+10\log \dfrac{I_1}{I_0}=93\Leftrightarrow 10\log \dfrac{I_1}{I_0}=93-10\log 7\approx 84{,}5.
   \] 
   Lúc sau mức cường độ  âm là $4I_1$.
    \[L=10\log \dfrac{4I_1}{I_0}=10\log 4+10\log \dfrac{I_1}{I_0}=93\Leftrightarrow 10\log \dfrac{I_1}{I_0}=10\log 4+84{,}5\approx 90{,}5\mathrm{dB}.
    \]
   \itemch Sai. Vì  nếu chỉ một bàn phím thì $L=10\log \dfrac{I}{I_0}=10\log \dfrac{10^{-5}}{10^{-12}}=70\mathrm{dB}$.\\
   Cả hai bàn phím cùng gõ thì  \[L_2=10\log \dfrac{2I}{I_0}=10\log 2+10\log \dfrac{I}{I_0}=10\log 2+70\approx 73\mathrm{dB}.\]
  \end{itemchoice}
 } 
\end{ex}

% nguồn: Ok_dungsai_lop11_Lan1
\begin{ex}%[1D6C4-6]%[ma:72351e7d]
 Cô Minh lần đầu gởi vào ngân hàng $100$ triệu đồng với kỳ hạn $3$ tháng,  lãi suất $2\%$ một quý theo hình thức lãi kép.
  \choiceTF
  {\True Sau $6$ tháng cô Minh  có tổng số tiền là $104{,}04$ triệu}
   {Để số tiền nhận được là $150$ triệu thì cô Minh phải gởi ngân hàng $18$ quý}
  {Sau đúng $6$ tháng cô Minh gởi thêm $100$ triệu đồng với kỳ hạn và lãi suất  như trước đó. Tổng số tiền cô Minh nhận được $1$ năm sau khi gởi thêm tiền gần nhất là $216$ triệu}
  {Để nhận được số tiền $200$ triệu trong 30 tháng với lãi suất như trên thì ban đầu cô Minh phải gởi ít  nhất $164$ triệu}
  \loigiai{
  \begin{itemchoice}
   \itemch Đúng. $3$ tháng $=1$ quý, nên $6$ tháng  $=2$ quý.\\
   Sau $6$ tháng cô Minh có tổng  số tiền là $100(1+2\%)^2=104{,}04$ triệu.
   \itemch Sai. Ta có $100(1+2\%)^n=150\Leftrightarrow n=\log _{\left(1+2\% \right)} 1{,}5\approx  20{,}475$.
   \itemch Sai. Cô Minh gởi thêm $100$ triệu đồng nên tổng số  tiền gởi là $204{,}04$ triệu.\\
   Số tiền sau $1$ năm nữa là $204{,}04(1+2\%)^4 \approx  220$ triệu.
   \itemch Sai. Để nhận  được số tiền $200$ triệu trong $30$ tháng với lãi suất như trên thì ban đầu cô Minh phải gởi ít nhất số tiền $x(1+2\%)^{10}=200\Leftrightarrow x\approx 164{,}06966$ triệu.
  \end{itemchoice}
 } 
\end{ex}

% nguồn: Ok_dungsai_lop11_Lan1
\begin{ex}%[1D6C4-6]%[ma:42ad818f]
 Trong Vật lí, sự phân rã của các chất phóng xạ được cho bởi công thức \break$m(t)=m_0\cdot \left(\dfrac{1}{2} \right)^{\frac{t}{T}}$; trong đó $m_0$ là khối lượng chất phóng xạ bạn đầu, $m(t)$ là khối lượng chất phóng xạ tại thời điểm $t$ và $T$ là chu kì bán rã. Hạt nhân Poloni là  chất phóng xạ $\alpha$ có chu kì bán rã là $138$ ngày. Giả sử lúc đầu có $100$ gam Poloni.
 \choiceTF 
  {\True Sau $138$ ngày thì khối lượng Poloni còn lại là $50$ (gam)} 
  {\True Khối lượng  Poloni còn lại sau $30$ ngày nhiều hơn $85$ (gam)}
  {\True Kể từ ngày thứ $55$ trở đi thì khối lượng Poloni còn lại ít hơn $75$  (gam)}
  {Kể từ ngày thứ $117$ trở đi thì khối lượng Poloni mất đi nhiều hơn $80\%$ so với khối lượng Poloni còn  lại}
  \loigiai{
  \begin{itemchoice}
    \itemch Đúng. Sau $138$ ngày thì khối lượng Poloni còn lại là \[m(138)=100\cdot \left(\dfrac{1}{2} \right)^{\frac{138}{138}}=50\text{ (gam).}\]
   \itemch Khối lượng Poloni còn  lại sau $30$ ngày là $m(30)=100\cdot\left(\dfrac{1}{2} \right)^{\frac{30}{138}} \approx 86,01 $ (gam).
   \itemch 
   \begin{eqnarray*}
    100\cdot\left(\dfrac{1}{2} \right)^{\frac{t}{138}} < 75&\Leftrightarrow& \left(\dfrac{1}{2} \right)^{\frac{t}{138}} < 0{,}75\\ 
    &\Leftrightarrow& \dfrac{t}{138} > \log _{\frac{1}{2}} 0{,}75\\ 
    &\Leftrightarrow& t  > 138\log _{\frac{1}{2}} 0{,}75\approx 57{,}28.
   \end{eqnarray*}
   Kể từ ngày thứ $58$ trở đi thì khối lượng Poloni còn lại ít hơn  $75$ (gam).
   \itemch  Khối lượng Poloni mất đi nhiều hơn $80\%$ so với khối lượng Poloni còn lại.\\
   Giả sử lúc đầu có $100$ gam Poloni, khối lượng Poloni còn lại là $x$  (gam).\\
   Khối lượng Poloni  mất đi bằng $100-x$ (gam).\\
   Suy ra $100-x > 80\% x\Rightarrow x+80\% x  < 100\Rightarrow 1{,}8x < 100\Rightarrow x < \dfrac{100}{1{,}8}=\dfrac{500}{9}$
   \[m(t)=100\cdot\left(\dfrac{1}{2} \right)^{\frac{t}{138}} < \dfrac{500}{9} \Rightarrow \left(\dfrac{1}{2} \right)^{\frac{t}{138}} < \dfrac{5}{9} \Rightarrow \dfrac{t}{138} > \log _{\frac{1}{2}} \dfrac{5}{9} \Rightarrow  t > 138\cdot\log _{\frac{1}{2}} \dfrac{5}{9} \approx 117{,}02.
    \]
    Vậy phải kể từ ngày thứ $118$ trở đi thì khối lượng Poloni mất đi nhiều hơn $80\%$ so với khối lượng Poloni còn lại.
  \end{itemchoice}
  }
\end{ex}

% nguồn: Ok_dungsai_lop11_Lan1
\begin{ex}%[1D6C4-6]%[ma:192ded33]
 Một người gửi tiết kiệm vào ngân hàng $1$ tỷ đồng. Các khẳng định  sau đây đúng hay sai?
 \choiceTF 
  {$8$ tháng sau người đó lấy về tất cả số tiền cả gốc và lãi là $1\,020\,175\,878$  đồng}
  {\True $2$ năm sau thì người  đó thu được số tiền số tiền cả gốc và lãi là $1\,127\,159\,776$ đồng}
  {\True Kể từ lúc gửi sau mỗi tháng vào ngày ngân hàng tính lãi người đó rút $10$ triệu đồng để chi tiêu thì  đến 139 tháng người đó rút hết tiền trong tài khoản}
  {\True Chưa đầy một năm thì lãi suất tăng lên thành $1{,}15\% /$ tháng. Tiếp theo, sáu tháng sau lãi suất chỉ còn $0{,}9\% /$ tháng. Người đó  tiếp tục gửi thêm một số tháng nữa rồi rút cả vốn lẫn lãi được $1\,143\,816\,503$ đồng. Vậy người đó gửi tiền vào ngân hàng với tổng thời gian là $16$ tháng}
  \loigiai{
  \begin{itemchoice}
   \itemch Sai. Đặt $T_0=1\,000\,000\,000$. Theo công thức lãi kép, số  tiền người đó nhận được sau $8$ tháng là $T_8=T_0 (1+r)^8$ thay $T_0=1\,000\,000\,000$ và $r=0{,}5\%$ ta được $T_8=1\,040\,707\,044$.
   \itemch Đúng. Theo công thức  lãi kép
   $T_{24}=T_0 \left(1+r\right)^{24}$ thay $T_0=1\,000\,000\,000$ và $r=0{,}5\%$ ta được  $T_{24}=1\,127\,159\,776$.
    \itemch Đúng. Ta có số tiền người đó gửi ban đầu là $T_0=1\,000$ triệu đồng, lãi suất hàng tháng $m=0{,}005$; số tiền người đó rút ra hàng tháng là $r=10$ triệu đồng.
   \begin{itemize}
     \item Sau tháng thứ nhất người đó thu được số tiền là $T_1=T_0 \left(1+m\right)$.
    \item Đầu tháng thứ hai người đó có  số tiền là $T_0 \left(1+m\right)-r$.
    \item Cuối  tháng thứ hai người đó có số tiền là \[T_2=\left(T_0 \left(1+m\right)-r\right)\left(1+m\right)=T_0 \left(1+m\right)^2-r\left(1+m\right).\]
    \item Đầu tháng  thứ ba người đó có số tiền là $T_0 \left(1+m\right)^2-r\left(1+m\right)-r$.
    \item Cuối tháng thứ ba người đó có  số tiền là \[T_3=T_0 \left(1+m\right)^3-r\left(1+m\right)^2-r\left(1+m\right).\]
    \item Cứ như thế số tiền người đó có cuối tháng thứ  $n$ là
    \begin{eqnarray*}
     T_n&=&T_0 \left(1+m\right)^n-\left[r\left(1+m\right)^{n-1}+r\left(1+m\right)^{n-2}+\cdots+r\left(1+m\right)\right]\\ 
      &=&T_0 \left(1+m\right)^n-r\cdot\dfrac{\left(1+m\right)^n-\left(1+m\right)}{m}.
    \end{eqnarray*}
    \item Người đó rút hết tiền trong  tài khoản khi
    \begin{eqnarray*}
     T_n-r\le 0&\Leftrightarrow& T_n \le 10\\ 
     &\Leftrightarrow& T_0  \left(1+m\right)^n-r\cdot\dfrac{\left(1+m\right)^n-\left(1+m\right)}{m} \le 10.
    \end{eqnarray*}
    Thay số ta được \[1\,000\cdot 1{,}005^n-10\cdot\dfrac{1{,}005^n-1{,}005}{0{,}005}  \le 10\Leftrightarrow 1{,}005^n \ge 2\Leftrightarrow n\ge 138{,}975.\]
   \end{itemize}
   \itemch Đúng. Gọi $m$  là số tháng người đó gửi tiền vào ngân hàng với lãi suất $0{,}5\% /$ tháng $(m\in \mathbb{N},1\le m\le 11)$.\\
   Khi đó, số tháng  người đó gửi tiền vào ngân hàng với lãi suất $1{,}15\% /$ tháng là $12-m$.\\
   Gọi $n$ là số tháng người đó gửi tiền vào ngân hàng  với lãi suất $0{,}9\% /$ tháng $\left(n\in \mathbb{N},1\le n\right)$.\\
   Tổng số tiền cả vốn lẫn lãi người đó  có được cho đến khi rút là \[T=T_0 \cdot (1{,}005)^m \cdot (1{,}0115)^{12-m} \cdot (1{,}009)^n,\text{ với } T_0=1\,000\,000\,000.\]
   Theo đề bài  thì $T=1\,143\,816\,503$.
   \begin{enumerate}[\bf Bước 1.]
    \item Với  $1\le m\le 11$, ta có đánh giá
    \begin{eqnarray*}
      & &T_0 \cdot (1{,}005)^{11} \cdot (1{,}0115)\cdot (1{,}009)^n \le T\le T_0 \cdot (1{,}005)\cdot (1{,}0115)^{11} \cdot (1{,}009)^n\\
     &\Leftrightarrow& \log _{1{,}009} \left(\dfrac{T}{T_0 \cdot (1{,}005)\cdot (1{,}0115)^{11}} \right)\le n\le  \log _{1{,}009} \left(\dfrac{T}{T_0 \cdot (1{,}005)^{11} \cdot (1{,}0115)} \right)\\
     &\Rightarrow& 1\le  n\le 7.
    \end{eqnarray*}
    \item Với $1\le  n\le 7$, ta đánh giá được
    \begin{eqnarray*}
     & &T_0 \cdot (1{,}005)^m \cdot (1{,}0115)^{12-m} \cdot (1{,}009)^1 \le T\le T_0 \cdot  (1{,}005)^m \cdot (1{,}0115)^{12-m} \cdot (1{,}009)^7\\
     &\Leftrightarrow& \dfrac{T}{T_0 \cdot (1{,}0115)^{12} \cdot (1{,}009)^7} \le \left(\dfrac{1{,}005}{1{,}0115} \right)^m \le \dfrac{T}{T_0 \cdot (1{,}0115)^{12}  \cdot (1{,}009)^1}\\
     &\Rightarrow& 2\le  m\le 10.
    \end{eqnarray*}
   \end{enumerate}
   Lặp lại các bước $1{,}2$ với các giá trị $m, n$ mới tìm được. Sau một số lần ta tìm được  $m=6$, $n=4$.\\
   Thử lại $m=6$,  $n=4$ thỏa mãn bài toán.\\
   Vậy người đó gửi tiền vào ngân  hàng với tổng thời gian là $16$ tháng.
  \end{itemchoice}
  }
\end{ex}

% nguồn: Ok_dungsai_lop11_Lan1
\begin{ex}%[1D6C4-6]%[ma:163e43f1]
 Sự tăng trưởng của một loại vi khuẩn tuân theo công thức $S=A\cdot \mathrm{e}^{rt}$, trong đó $A$ là số lượng vi khuẩn ban đầu, $r$ là tỉ lệ tăng trưởng, $t$ là thời gian tăng trưởng. Biết rằng số lượng vi khuẩn ban đầu là $100$ con và  sau $5$ giờ có $300$ con.
  \choiceTF
   {\True Phương trình thể hiện tỷ lệ tăng trưởng là $300=100\cdot \mathrm{e}^{5r}$}
  {Tỉ lệ tăng trưởng của loại  vi khuẩn này là $r=\ln \dfrac{3}{5}$ mỗi giờ}
  {\True  Sau $10$ giờ, từ $100$ con vi khuẩn sẽ có $100\cdot \mathrm{e}^{10\cdot\frac{\ln 3}{5}}=900$ con}
  {Phải cần ít nhất $20$ giờ để số con vi khuẩn lớn hơn $10\,000$  con}
  \loigiai{
  \begin{itemchoice}
    \itemch Đúng. Từ giả thiết ta có $300=100\cdot \mathrm{e}^{5r}$.
   \itemch Sai.  Ta có $300=100\cdot \mathrm{e}^{5r} \Leftrightarrow r=\dfrac{\ln 300-\ln 100}{5}=\dfrac{\ln 3}{5}$.\\
   Tức  tỉ lệ tăng trưởng của loại vi khuẩn này là $r=\dfrac{\ln3}{5}$ mỗi giờ.
   \itemch Đúng. Sau  $10$ giờ, từ $100$ con vi khuẩn sẽ có $100\cdot\mathrm{e}^{10\cdot\frac{\ln 3}{5}}=900$ con.
   \itemch Sai. Để  số con số vi khuẩn lớn hơn $10\,000$ con thì \[100\cdot \mathrm{e}^{t\cdot\frac{\ln 3}{5}} > 10\,000\Leftrightarrow \mathrm{e}^{t\cdot\frac{\ln 3}{5}} > 100\Leftrightarrow t\cdot\dfrac{\ln 3}{5} > \ln 100\Leftrightarrow t > \dfrac{5\ln 100}{\ln 3} \approx 20{,}96.\]
  \end{itemchoice}
  }
\end{ex}

% nguồn: Ok_dungsai_lop11_Lan1
\begin{ex}%[1D6C4-6]%[ma:12aea4a7]
 Dân số nước ta năm $2\,023$ ước tính là $A=100{,}3$ triệu người. Giả sử tỉ lệ tăng dân số  hằng năm của nước ta là $r=0{,}84\%$. Biết rằng sau $t$ năm, dân số Việt Nam được tính theo công thức $S=A\cdot \mathrm{e}^{rt}$ triệu người.
 \choiceTF 
  {\True Sau $1$ năm  nữa dân số Việt Nam đạt $101{,}1$ triệu người}
  {Đến  năm $2\,030$, dân số Việt Nam ước đạt $120$ triệu người}
  {Người ta ước tính rằng, đến năm $2\,035$. Mức sinh của Việt Nam có xu hướng giảm, tỉ lệ  tăng dân số hằng năm chỉ còn khoảng $r=0{,}4\%$. Dân số Việt Nam vào năm $2\,040$ là hơn $120$ triệu người}
  {Dân  số nước ta vượt $110$ triệu người trong vòng $10$ năm nữa}
 \loigiai{ 
  \begin{itemchoice}
   \itemch Đúng. Sau $1$ năm nữa dân số Việt Nam  ước tính là \[100{,}3\cdot \mathrm{e}^{0{,}0084\cdot 1}=101{,}1\text{ (triệu người).}\]
   \itemch Sai. Đến năm $2030$, tức là sau $7$ năm dân số Việt  Nam ước tính là \[100{,}3\cdot \mathrm{e}^{0{,}0084\cdot 7}=106{,}4\text{ (triệu người).}\]
   \itemch Sai. Với  tỉ lệ tăng dân số hằng năm là $r=0{,}84\%$, ước tính dân số Việt Nam vào năm 2035 là $100{,}3\cdot \mathrm{e}^{0{,}0084\cdot 12}=110{,}9$ triệu người.
   Từ $2\,035$, tỉ lệ  tăng dân số hằng năm chỉ còn khoảng $r=0{,}4\%$, nên tính đến $2\,040$, dân số Việt Nam khoảng $110{,}9\cdot \mathrm{e}^{0{,}004\cdot 5}=113{,}1$ triệu người.
   \itemch Sai. Ta có $100{,}3\cdot \mathrm{e}^{0{,}0084\cdot t} > 110\Leftrightarrow \mathrm{e}^{0{,}0084\cdot t} > \dfrac{110}{100{,}3} \Leftrightarrow 0{,}0084t >  \ln \dfrac{110}{100{,}3} \Leftrightarrow t > 10{,}9$.\\
   Vậy sau $11$ năm nữa tính từ mốc $2\,023$,  tức là từ năm $2\,034$ trở đi, dân số Việt Nam ước tính vượt quá $110$ triệu người.
  \end{itemchoice}
  }
\end{ex}


%%%%%%%%%%%%%%%% TRẢ LỜI NGẮN %%%%%%%%%%%%%%%%

% nguồn: Ok_TraLoiNgan_Lop11_Lan1
\begin{ex}%[1D6H4-6]%[ma:611a80cd]
	Một ngân hàng X quy định về số  tiền nhận được của khách hàng sau $n$ năm gửi tiền vào ngân hàng tuân theo công thức $P(n)=A \left(1+8\%\right)^n$, trong đó $A$ là số tiền gửi ban đầu của khách hàng. Hỏi số tiền $A$ (triệu đồng) ít nhất mà khách hàng phải gửi là bao nhiêu để sau $3$ năm khác hàng đó nhận được lớn hơn $850$ triệu đồng (kết quả làm tròn đến hàng đơn vị).
	 \shortans{675}
	\loigiai{Để  sau $3$ năm khách hàng nhận được lớn hơn $850$ triệu đồng thì số tiền gửi $A$ thỏa
	 \allowdisplaybreaks
	\begin{align*}
		P(n)>850&\Leftrightarrow A \left(1+8\%\right)^3>850\\ 
		&\Leftrightarrow A>\dfrac{850}{\left(1+8\%\right)^3} \approx  675 .
	\end{align*}
	Vậy để nhận được  lớn hơn $850$ triệu đồng sau $3$ năm, khách hàng đó phải gửi vào ít nhất $675$ triệu đồng.
	} 
\end{ex}

% nguồn: Ok_TraLoiNgan_Lop11_Lan1
\begin{ex}%[1D6H4-6]%[ma:83f1b786]
	Áp suất khí quyển $p$ (tính bằng kilôpascan, viết tắt là $kPa$) ở độ cao $h$ (so với mực nước biển, tính bằng $km$) được  tính theo công thức sau: $\ln\left(\dfrac{p}{100}\right)=-\dfrac{h}{7}$. Ở độ cao $12km$ thì áp suất khí quyển là bao nhiêu kPa? (Làm tròn kết quà đến hàng đơn vị).\shortans{18}
	\loigiai{ 
		Thay $h=12$  vào $\ln\left(\dfrac{p}{100}\right)=-\dfrac{h}{7}$, ta có
		 $$\ln\left(\dfrac{p}{100}\right)=-\dfrac{12}{7}\Leftrightarrow \dfrac{p}{100}=e^{-\dfrac{12}{7}} \Leftrightarrow p=100\cdot e^{-\dfrac{12}{7}}\approx 18.$$
		Vậy ở độ cao $12km$ thì  áp suất khí quyển khoảng $18$ kPa.
	 }
\end{ex}

% nguồn: Ok_TraLoiNgan_Lop11_Lan1
\begin{ex}%[1D6H4-6]%[ma:33de1069]
	Sự tăng trưởng của một loài vi khuẩn  được tính theo công thức $f(t)=A\cdot e^{rt}$, trong đó $A$ là số lượng vi khuẩn ban đầu, $r$ là tỷ lệ tăng trưởng $(r>0)$, $t$ (tính theo giờ) là thời gian tăng trưởng. Biết số vi khuẩn ban đầu có $1000$ con và sau $10$ giờ là $5000$ con. Hỏi sau ít nhất bao nhiêu giờ thì số lượng vi khuẩn tăng gấp $10$ lần (làm tròn kết quả đến hàng phần chục)?
	\shortans{14,3} 
	\loigiai{ 
		Ta có $5000=1000 e^{10 r} \Leftrightarrow e^{10 r}=5 \Leftrightarrow r=\dfrac{\ln  5}{10}$.\\
		Gọi $t$ (giờ) là thời  gian để số lượng vi khuẩn tăng gấp $10$ lần.\\
		 Ta có
		\begin{eqnarray*}
			& & 10000=1000 e^{r t}\\ 
			&\Leftrightarrow & e^{r t}=10\\ 
			&\Leftrightarrow & r  t=\ln 10\\
			 &\Leftrightarrow & t=\dfrac{\ln 10}{r}\\
			&\Leftrightarrow & t=\dfrac{\ln 10}{\dfrac{\ln  5}{10}}\\
			&\Leftrightarrow &  t=10 \log _5 10\\
			&\Leftrightarrow & t\approx  14,3.\\
		\end{eqnarray*}
		Vậy sau khoảng  $14,3$ giờ thì số lượng vi khuẩn sẽ tăng gấp $10$ lần.
	} 
\end{ex}

% nguồn: Ok_TraLoiNgan_Lop11_Lan2
\begin{ex}%[1D6H4-6]%[ma:53590d07]
	Ông Năm gửi tiết kiệm $200$ triệu đồng với lãi suất $9{,}1\% /$ năm. Tính tổng số tiền vốn và lãi mà ông Năm nhận được sau một năm nếu tiền lãi được tính theo thể thức lãi kép với kì  hạn $3$ tháng (đơn vị triệu đồng, kết quả làm tròn đến hàng đơn vị).
	 \shortans{219}
	\loigiai{ 
		Ta có $A=200$; $r=14\%$ và  $n=\dfrac{12}{3}=4$.\\
		Vậy tổng số tiền vốn và lãi mà ông Năm nhận được sau một năm nếu tiền lãi được tính theo thể thức lãi kép với kì  hạn $3$ tháng là \[T=A\left(1+\dfrac{r}{n} \right)^n=200\cdot \left(1+\dfrac{9{,}1\%}{4} \right)^4\approx 218{,}8\text{ (triệu đồng).}\]
	 }
\end{ex}

% nguồn: Ok_TraLoiNgan_Lop11_Lan2
\begin{ex}%[1D6H4-6]%[ma:55b75683]
	Người ta sử dụng công thức $S=A\cdot \mathrm{e}^{n\cdot r}$ để dự báo dân số của một quốc gia, trong đó $A$ là số dân của năm lấy làm mốc tính, $S$ là số dân sau $n$ năm và $r$ là tî lệ gia  tăng dân số hàng năm. Biết rằng năm $2\,001$, dân số của Việt Nam là $78\,685\,800$ người. Giả sử tỉ lệ tăng dân số hàng năm không đổi là $1{,}2\%$. Hãy tính xem dân số nước ta đạt $110$ triệu người vào năm nào?
	\shortans{2029} 
	 \loigiai{
		Theo công thức tăng  trưởng mũ $S=A\cdot \mathrm{e}^{n\cdot r}$, ta có
		\[110\,000\,000=78\,685\,800\cdot \mathrm{e}^{1{,}2\%  \cdot n} \Leftrightarrow n=\dfrac{1}{1{,}2\%} \ln \dfrac{110\,000\,000}{78\,685\,800} \sim eq 27{,}91.
		\] 
		Vậy sau $28$ năm thì  dân số Việt Nam đạt $110$ triệu người và đạt $110$ triệu người vào năm $2\,029$.
	} 
\end{ex}

% nguồn: Ok_TraLoiNgan_Lop11_Lan2
\begin{ex}%[1D6H4-6]%[ma:88f47422]
	Anh Nam tiết kiệm được $x$ triệu đồng và dùng tiền đó để mua một căn nhà nhưng thực tế giá căn nhà đó là $1{,}6x$ triệu đồng. Anh Nam quyết định gửi tiết kiệm vào ngân hàng với lãi suất $7\% $/ năm theo hình thức lãi kép và không rút tiền trước kỳ hạn. Hỏi sau ít nhất bao nhiêu năm anh Nam có đủ số tiền cần thiết mua căn nhà  đó? Giả định trong suốt thời gian gửi, lãi suất không đổi, anh Nam không rút tiền ra và giá bán căn nhà không thay đổi.
	\shortans{7} 
	\loigiai{ 
		Số tiền anh Nam  có sau $n$ năm là $T_n =x(1+0{,}07)^n $, $(n\in \mathbb{N}^*)$.\\
		Yêu cầu bài toán $\Leftrightarrow x(1+0{,}07)^n \ge 1{,}6x \Leftrightarrow 1{,}07^n  \ge 1{,}6\Leftrightarrow n\ge \log _{1{,}07} 1{,}6\approx 6{,}94$.\\
	Vậy sau ít nhất $7$ năm anh Nam có đủ số tiền cần thiết mua căn nhà đó.} 
\end{ex}

% nguồn: Ok_TraLoiNgan_Lop11_Lan2
\begin{ex}%[1D6H4-6]%[ma:829b0cea]
	Số lượng vi khuẩn  trong một phòng thí nghiệm $A$ được tính theo công thức $s(t)=s(0)\cdot 2^t$, trong đó $s(0)$ là số lượng vi khuẩn $A$ lúc ban đầu, $s(t)$ là số lượng vi khuẩn sau $t$ phút. Biết sau $3$ phút thì số lượng vi khuẩn $A$ là 625 nghìn con. Tính thời gian kể từ lúc ban đầu, số lượng loại vi khuẩn $A$ là $20$ triệu con?
	 \shortans{8}
	\loigiai{ 
		Ta có $s(3)=625$ nghìn  con $\Rightarrow s(0)\cdot 2^3=625\Rightarrow s(0)=\dfrac{625}{8}$ nghìn con.\\
		 Để số lượng vi khuẩn là $20$ triệu con thì \[s(0)\cdot 2^t=20000\Rightarrow 2^t=20000:\dfrac{625}{8}\Rightarrow 2^t=2^8\Rightarrow t=8.\]
		Vậy thời gian để số vi khuẩn đạt  $20$ triệu con là $8$ phút.
	} 
\end{ex}

% nguồn: Ok_TraLoiNgan_Lop11_Lan2
\begin{ex}%[1D6H4-6]%[ma:8812d71d]
	Kết quả thống kê cho biết ở thời điểm năm $2013$ dân số Việt Nam là $90$ triệu người, tốc độ tăng dân số là $1{,}1\% /$năm. Nếu mức tăng dân số ổn định như vậy thì dân số  Việt Nam sau $t$ năm kể từ năm $2013$ được tính bởi công thức $P(t)=90\cdot(1+1{,}1\%)^t $. Hỏi đến năm $2077$ dân số Việt Nam là khoảng bao nhiêu triệu người? (kết quả làm tròn đến hàng đơn vị)
	\shortans{181} 
	\loigiai{ 
		Ta có Tính từ năm $2013$ đến  năm $2077$ là $64$ năm.\\
		Áp dụng công thức $P(t)=90\cdot(1+1{,}1\%)^t $ thì đến năm $2077$ dân số  nước ta là
		\[P=90\cdot(1+1{,}1\%)^{64} \approx  181{,}266\, (\text{triệu người}).\]
	 }
\end{ex}

% nguồn: Ok_TraLoiNgan_Lop11_Lan2
\begin{ex}%[1D6H4-6]%[ma:bee8db6c]
	Biết rằng năm $2001$, dân số Việt Nam là $78\,685\,800$ người và tỉ lệ tăng dân số năm đó là $1{,}7\% $. Cho biết sự tăng dân số được ước tính  theo công thức $S=A\cdot \mathrm{e}^{Nr} $. Nếu dân số vẫn tăng với tỉ lệ như vậy thì đến năm nào dân số nước ta ở mức $120$ triệu người?
	 \shortans{2,03}
	 \loigiai{
		Ta có  Nếu dân số vẫn tăng với tỉ lệ $1{,}7\% $ thì sau $N$ năm dân số Việt Nam là $120\,000\,000$ người.\\
		Ta có  $S=A\cdot \mathrm{e}^{Nr} \Leftrightarrow 120\,000\,000=78\,685\,800\cdot \mathrm{e}^{1{,}7\% N} \Rightarrow N\approx 25$.\\
	Vậy đến năm $2001+25=2026$ thì dân số nước ta ở mức $120$  triệu người}
\end{ex}

% nguồn: Ok_TraLoiNgan_Lop11_Lan2
\begin{ex}%[1D6H4-6]%[ma:e103dd86]
	Bác An gửi tiết kiệm ngân hàng $100\cdot10^6 $ đồng kì hạn $12$ tháng, với lãi suất không  đổi là $6\% $ một năm. Khi đó sau $n$ năm gửi thì tổng số tiền bác An thu được cho bởi công thức sau: $A=100\cdot10^6 (1+0{,}06)^n $ đồng. Hỏi sau ít nhất bao nhiêu năm, tổng số tiền bác An thu được là không dưới $150\cdot10^6 $đồng?
	\shortans{7} 
	 \loigiai{
		Theo giả  thiết ta có $100\cdot 10^6 (1+0{,}06)^n \ge 150\cdot10^6 \Rightarrow n\ge \log _{1{,}06} 1{,}5\approx 6{,}96$.\\
	Vậy sau ít nhất $7$ năm thì bác An nhận được số tiền ít nhất là $150\cdot10^6 $  đồng.}
\end{ex}

% nguồn: Ok_TraLoiNgan_Lop11_Lan2
\begin{ex}%[1D6H4-6]%[ma:084f1daf]
	Giá trị còn lại của một chiếc xe mua mới theo thời gian $t$ được xác định bởi công thức: $V(t)=1{,}5\mathrm{e}^{-0{,}15t} $, trong đó $V(t)$ được tính bằng tỷ đồng và $t$ tính bằng năm. Sau ít nhất bao nhiêu năm kể  từ thời điểm mua xe giá trị chiếc xe đó còn lại dưới $500$ triệu đồng?
	 \shortans{8}
	 \loigiai{
		Ta có $V(t)=1{,}5 \mathrm{e}^{-0{,}15t}  $.\\
		Nên khi giá xe  còn lại dưới $500$ triệu đồng thì ta có
		\[0{,}5\ge  1{,}5 \mathrm{e}^{-0{,}15t} \Leftrightarrow t\ge 7{,}32.\]
	Vậy sau ít nhất $8$ năm kể từ  thời điểm mua xe giá trị chiếc xe đó sẽ còn lại dưới $500$ triệu đồng.}
\end{ex}

% nguồn: Ok_TraLoiNgan_Lop11_Lan2
\begin{ex}%[1D6H4-6]%[ma:bd178d75]
	Một người gửi tiết kiệm vào ngân hàng với lãi suất $7{,}5\%$/năm. Biết rằng nếu không  rút tiền ra khỏi ngân hàng thì cứ sau mỗi năm số tiền lãi sẽ được nhập vào vốn để tính lãi cho năm tiếp theo. Tính sau thời gian ngắn nhất (theo năm) để số tiền người đó thu được (cả số tiền gửi ban đầu và lãi) gấp đôi số tiền đã gửi ban đầu, giả định trong khoảng thời gian này lãi suất không thay đổi và người đó không rút tiền ra?
	\shortans{10} 
	\loigiai{ 
		Giả sử số tiền người đó gửi vào ngân  hàng là $A$. \\
		Sau $n$ năm số tiền người đó nhận  được là $2A$.\\
		Áp dụng công thức $S=A\mathrm{e}^{0{,}075t}$, ta  có
		\begin{eqnarray*}
			&&2A\le  A\cdot \mathrm{e}^{0{,}075t}\\
			&\Leftrightarrow& 0{,}075t\ge \ln  2\\
			&\Leftrightarrow& t\ge \dfrac{\ln 2}{0{,}075}\approx 9{,}24. 
		\end{eqnarray*}
		Vậy  người đó phải gửi ít nhất $10$ năm thì số tiền thu được gấp đôi số tiền ban đầu.
	 }
\end{ex}

% nguồn: Ok_TraLoiNgan_Lop11_Lan2
\begin{ex}%[1D6H4-6]%[ma:fd6f50f1]
	Số lượng loại vi khuẩn $A$ trong một phòng thí nghiệm được tính theo công thức $s(t)=s(0)\cdot2^t  $, trong đó $s(0)$ là số lượng vi khuẩn $A$ lúc ban đầu, $s(t)$ là số lượng vi khuẩn $A$ có sau $t$ phút. Biết sau $3$ phút thì số vi khuẩn $A$ là $625$ nghìn con. Hỏi sau bao lâu kể từ lúc ban đầu, số lượng loại vi khuẩn $A$ là $20$ triệu con.
	\shortans{8} 
	\loigiai{ 
		Ta  có $s(3)=625\,000\Leftrightarrow s(0)\cdot2^3 =625\,000\Leftrightarrow s(0)=78\,125$.\\
		Số lượng loại vi khuẩn $A$ là  $20$ triệu con khi
		\[s(t)=20\,000\,000\Leftrightarrow s(0)\cdot2^t =20\,000\,000\Leftrightarrow 2^t =\dfrac{20\,000\,000}{s(0)} =\dfrac{20\,000\,000}{78\,125} =256\Leftrightarrow t=8.\] 
	Vậy, sau $8$ phút thì số lượng vi khuẩn $A$ là $20$  triệu con.}
\end{ex}

% nguồn: Ok_TraLoiNgan_Lop11_Lan1
\begin{ex}%[1D6V4-6]%[ma:2f608aa3]
	Anh Hưng gửi tiết kiệm khoản tiền 700 triệu đồng vào một ngân hàng  với lãi suất $7\%$ / năm theo hình thức lãi kép kì hạn 12 tháng. Tính thời gian tối thiểu gửi tiết kiệm để anh Hưng thu được ít nhất 1 tỉ đồng (cả vốn lẫn lãi). Cho biết công thức lãi kép là $T=A\cdot (1+r)^n$, trong đó $A$ là tiền vốn, $T$ là tiền vốn và lãi nhận được sau $n$ năm, $r$ là lãi suất/năm.
	 \shortans{6}
	\loigiai{ 
		Ta có: $T\ge 1000\Leftrightarrow 700(1+7\%)^n \ge  1000\Leftrightarrow 1,07^n \ge \dfrac{10}{7}$
		\[\Leftrightarrow  n\ge \log _{1,07} \left(\dfrac{10}{7} \right)\approx 5,27 \; (do\; 1,07 > 1).\;
		 \]
		Vậy thời gian gửi tiết kiệm phải ít nhất 6 năm thì  anh Hưng mới thu được ít nhât 1 tỉ đồng.
	} 
\end{ex}

% nguồn: Ok_TraLoiNgan_Lop11_Lan1
\begin{ex}%[1D6V4-6]%[ma:ee004c61]
	Một người gửi tiền vào ngân hàng với lãi suất không thay đổi là $ 6\% $ trên năm. Biết rằng nếu không rút tiền ra khỏi ngân hàng thì cứ sau mỗi năm, số tiền lãi sẽ được nhập  vào vốn ban đầu (người ta gọi đó là lãi kép). Người đó định gửi tiền trong vòng $ 3 $ năm, sau đó rút $ 500 $ triệu đồng. Hỏi số tiền ít nhất người đó phải gửi vào ngân hàng (kết quả làm tròn đến hàng triệu) là bao nhiêu triệu đồng?
	 \shortans{420}
	\loigiai{ 
		Gọi $ X $ (triệu đồng)  là số tiền ban đầu người đó gửi vào ngân hàng.\\
		Sau $ 3 $ năm sẽ có  $ X\cdot (1+0{,}06)^3 $.\\
		Theo đề bài ta có $ X\cdot  (1+0{,}06)^3 \geq 500\Leftrightarrow X\geq 419{,}8$.\\
		Vậy người đó phải gửi vào ngân hàng số tiền ít nhất là $ 420 $ triệu  đồng.
	} 
\end{ex}

% nguồn: Ok_TraLoiNgan_Lop11_Lan1
\begin{ex}%[1D6V4-6]%[ma:beb847a0]
	Người ta dùng thuốc để khử khuẩn cho một thùng nước. Biết rằng nếu lúc đầu mỗi mililít nước chứa $P_{0}$ vi khuẩn thì sau $t$ giờ (kể từ khi cho thuốc vào thùng), số lượng vi khuẩn trong mỗi mililít nước là $P=P_0 \cdot 10^{-\alpha t}$, với $\alpha$ là một hằng số dương nào đó. Biết rằng ban đầu mỗi mililít nước có $9000$ vi khuẩn và sau $2$ giờ, số lượng vi khuẩn trong mỗi mililit nước là $6000$. Sau thời gian bao lâu thì số lượng vi khuẩn trong  mỗi mililit nước trong thùng ít hơn hoặc bằng $1000$ (làm tròn đến hàng phần mười)?
	 \shortans{10,8}
	\loigiai{ 
		Ta có $6000=9000\cdot 10^{-2\alpha} \Rightarrow \alpha=-\dfrac{1}{2}  \log \dfrac{2}{3}=\dfrac{1}{2} \log \dfrac{3}{2}$.\\
		 \allowdisplaybreaks
		\begin{eqnarray*}
			9000 \cdot  10^{-\alpha t} \leq 1000
			 &\Leftrightarrow & 10^{-\alpha t} \leq \dfrac{1}{9} \\
			&\Leftrightarrow &-\alpha t \leq \log \dfrac{1}{9} \\ 
			&\Leftrightarrow &t \geq-\dfrac{2}{\alpha} \log \dfrac{1}{3}=-\dfrac{2}{\dfrac{1}{2} \log  \dfrac{3}{2}} \cdot \log \dfrac{1}{3}=\dfrac{4 \log 3}{\log \dfrac{3}{2}} \approx 10{,}8 \ \text{(giờ).}
		\end{eqnarray*}
	} 
\end{ex}

% nguồn: Ok_TraLoiNgan_Lop11_Lan1
\begin{ex}%[1D6V4-6]%[ma:85033764]
	Sau một tháng thi công, công trình xây dựng lớp học từ thiện cho học sinh vùng cao đã thực hiện được một khối lượng công việc. Nếu tiếp tục với tiến độ như  vậy thì dự kiến sau đúng 23 tháng nữa công trình sẽ hoàn thành. Để sớm hoàn thành công trình và kịp thời đưa vào sử dụng, đơn vị xây dựng quyết định từ tháng thứ hai tăng $5\%$ khối lượng công việc so với tháng kề trước. Hỏi công trình sẽ hoàn thành ở tháng thứ mấy sau khi khởi công?
	 \shortans{17}
	\loigiai{ 
		Theo dự kiến, cần $24$ tháng để hoàn thành công trình. Vậy khối lượng  công việc trên một tháng theo dự tính là $\dfrac{1}{24}$.\\
		Khối lượng công việc của tháng thứ  2 là $T_2=\dfrac{1}{24}+0{,}05\cdot\dfrac{1}{24}=\dfrac{1}{24}(1+0{,}05)^1$.\\
		Khối lượng công việc của tháng thứ 3  là\\
		$T_3=\left(\dfrac{1}{24}+0{,}05\cdot\dfrac{1}{24}\right)+0{,}05\cdot\left(\dfrac{1}{24}+0{,}05\cdot\dfrac{1}{24}\right) =\dfrac{1}{24}\cdot (1+0{,}05)^2$.\\ 
		Như vậy khối  lượng công việc của tháng thứ $n$ là $T_n=\dfrac{1}{24}\cdot (1+0{,}05)^{n-1}$.\\
		Ta có \allowdisplaybreaks\begin{eqnarray*}
		&&\dfrac{1}{24}\cdot  (1+0{,}05)^{\circ}+\dfrac{1}{24}\cdot (1+0{,}05)^1+\cdots +\dfrac{1}{24}\cdot (1+0{,}05)^{n-1}=1 \\
		&\Leftrightarrow&\dfrac{1}{24}\cdot\dfrac{1-(1+0{,}05)^n}{1-(1+0{,}05)}=1\\ 
		 &\Leftrightarrow&(1+0{,}05)^n=\dfrac{11}{5}\Leftrightarrow n=\log_{1+0{,}05}\dfrac{11}{5}\approx 16{,}2.
	\end{eqnarray*}
	Vậy công trình sẽ hoàn thành ở tháng thứ $17$  từ khi khởi công.
 }
\end{ex}

% nguồn: Ok_TraLoiNgan_Lop11_Lan1
\begin{ex}%[1D6V4-6]%[ma:043e8cd1]
	Một người lần đầu gửi vào ngân hàng $400$ triệu đồng với kì hạn $3$ tháng, lãi suất $4 \%$ một quý. Biết rằng  nếu không rút tiền ra khỏi ngân hàng thì cứ sau mỗi quý số tiền lãi sẽ được nhập vào gốc để tính lãi cho quý tiếp theo. Sau đúng $6$ tháng, người đó gửi thêm $100$ triệu đồng với kỳ hạn và lãi suất như trước đó. Tổng số tiền người đó nhận được $1$ năm sau khi gửi tiền (cả vốn lẫn lãi, đơn vị triệu đồng)?
	\shortans{576} 
	\loigiai{ 
		Số tiền lãi  và gốc sau $6$ tháng là $400(1+4\%)^2$ triệu đồng.\\
		Số tiền gốc  sau khi hết $6$ tháng thên $100$ triệu đồng là $400(1+4\%)^2+100$.\\
		Số tiền cả lãi  lẫn gốc sau $1$ năm là $[400(1+4\%)^2+100](1+4\%)^2\approx 576$ triệu đồng.
	 }
\end{ex}

% nguồn: Ok_TraLoiNgan_Lop11_Lan1
\begin{ex}%[1D6V4-6]%[ma:296d2a41]
	Một nhóm các chuyên gia y tế đang nghiên cứu và thử nghiệm độ chính xác của một xét nghiệm COVID-19. Giả sử cứ sau $n$ lần thử nghiệm và điều chỉnh thì tỉ lệ chính  xác của bộ xét nghiệm đó tuân theo công thức $S(n)=\dfrac{1}{1+2020 \cdot 10^{-0{,}01n}}$. Hỏi phải tiến hành ít nhất bao nhiêu lần thử nghiệm và điều chỉnh bộ xét nghiệm để đảm bảo độ chính xác của bộ xét nghiệm đo trên $90\%$?
	 \shortans{}
	\loigiai{Để đảm bảo độ chính xác của bộ xét  nghiệm đo trên $90\%$ thì số lần thử nghiệm $n$ thỏa
	\allowdisplaybreaks 
	\begin{align*}
		S(n)>90\% &\Leftrightarrow\dfrac{1}{1+2020 \cdot  10^{-0{,}01n}}>90\%\\
		&\Leftrightarrow 1+2020 \cdot 10^{-0{,}01n} <  \dfrac{100}{90}\\
		&\Leftrightarrow 2020 \cdot 10^{-0{,}01n}  <\dfrac{1}{9}\\
		&\Leftrightarrow 10^{-0{,}01n}<\dfrac{1}{18180}\\ 
		&\Leftrightarrow -0{,}01n< \log  \dfrac{1}{18180}\\
		 &\Leftrightarrow n> -\dfrac{1}{0{,}01}\log \dfrac{1}{18180}.
	\end{align*}
	Số tự nhiên $n$ nhỏ nhất thỏa  $n> -\dfrac{1}{0{,}01}\log \dfrac{1}{18180}$ là $426$.\\
	Vậy phải tiến hành ít nhất $426$ lần thử nghiệm và điều chỉnh bộ xét nghiệm để đảm bảo độ chính  xác của bộ xét nghiệm đo trên $90\%$.
	 }
\end{ex}

% nguồn: Ok_TraLoiNgan_Lop11_Lan1
\begin{ex}%[1D6V4-6]%[ma:40441fa4]
	Anh Dũng đem gửi tiết kiệm số tiền là $400$ triệu đồng ở hai  loại kì hạn khác nhau. Anh gửi $250$ triệu đồng theo kì hạn $3$ tháng với lãi suất $x\%$ một quý, số tiền còn lại anh gửi theo kì hạn $1$ tháng với lãi suất $0{,}25\%$ một tháng. Biết rằng nếu không rút lãi thì số lãi được nhập vào số tiền gốc để tính lãi cho kì hạn tiếp theo. Sau $1$ năm, số tiền cả gốc lẫn lãi của anh là $416780000$ đồng. Tính $x$ (làm tròn đến hàng phần trăm).
	 \shortans{1,2}
	\loigiai{Giả sử anh Dũng gửi vào số tiền $A$ triệu đồng với lãi suất $x\%$ một quý. Khi đó số  tiền cả vốn lẫn lãi anh nhận được sau $n$ quý là $A(1+r\%)^n$.\\
	Ở loại kì hạn thứ nhất, anh Dũng gửi $250$ triệu đồng theo kì hạn $1$ quý ($3$ tháng) với lãi suất $x\%$ một quý. Do đó sau $1$ năm ($4$ quý) số tiền cả vốn lẫn  lãi anh thu được là $A_1=250(1+x\%)^4$.\\
	Ở loại kì  hạn thứ hai, anh Dũng gửi $150$ triệu đồng (số tiền còn lại) theo kì hạn $1$ tháng với lãi suất $0{,}25\%$ một tháng. Do đó sau $1$ năm ($12$ tháng) số tiền cả vốn lẫn lãi anh thu được là $A_2=150(1+0{,}25\%)^{12}$.\\
	Vì sau $1$ năm, số tiền cả  gốc lẫn lãi của anh là $416780000$ đồng nên
	\allowdisplaybreaks 
	\begin{align*}
		A_1+A_2=416780000&\Leftrightarrow 250 \cdot 10^6  \cdot (1+x\%)^4+150 \cdot 10^6 \cdot (1+0{,}25\%)^{12}=416780000\\
		&\Leftrightarrow  (1+x\%)^4= \dfrac{416780000-150 \cdot 10^6 \cdot (1+0{,}25\%)^{12}}{250 \cdot 10^6}\\
		&\Leftrightarrow  x\%=\sqrt[4]{\dfrac{416780000-150 \cdot 10^6 \cdot (1+0{,}25\%)^{12}}{250 \cdot 10^6}}-1\\
		&\Leftrightarrow x  \approx 1{,}2.
	\end{align*}
	Vậy $x \approx 1{,}2$. 
	 }
\end{ex}

% nguồn: Ok_TraLoiNgan_Lop11_Lan1
\begin{ex}%[1D6V4-6]%[ma:c8f7bf63]
	Công thức tính khối lượng còn lại của một chất phóng xạ từ khối  lượng ban đầu $m_0$ được cho bởi công thức:
	$$m(t)=m_0\left(\dfrac{1}{2}\right)^{\tfrac{t}{T}},$$ 
	trong đó $t$ là thời gian tính từ thời điểm ban đầu và $T$ là chu kì bán rã  của chất đó. Biết rằng chất phóng xạ polonium-210 có chu kì bán rã là $ 138 $ ngày. Từ khối lượng polonium-210 ban đầu $100 \mathrm{~g}$, sau bao lâu khối lượng còn lại là $50 \mathrm{~g}$?
	(Kết quả tính theo ngày  và làm tròn đến chữ số thập phân thứ hai).
	 \shortans{138}
	 \loigiai{
		Giải phương trình  $100\left(\dfrac{1}{2}\right)^{\tfrac{t}{138}}=50$, ta được $t=138$.\\
		Vậy  sau $ 138 $ ngày thì khối lượng polonium-210 còn $50 \mathrm{~g}$.
	} 
\end{ex}

% nguồn: Ok_TraLoiNgan_Lop11_Lan1
\begin{ex}%[1D6V4-6]%[ma:65436736]
	Dân số ở một địa phương được ước tính theo công thức $S=A\cdot e^{r.t}$, trong đó $A$ không đổi là dân số của năm 2023, $S$ là dân số sau $t$ năm, $r$ là tỉ lệ tăng dân số hằng năm. Hỏi đến năm nào thì dân số ở địa phương đó sẽ đạt gấp đôi dân số năm  2023? Biết $r=1,13\% /$ năm.
	\shortans{2085} 
	 \loigiai{
		Dân  số đạt gấp đôi nghĩa là $S=2A$, ta có
		\[2A=A\cdot e^{1,13\%.t} \Leftrightarrow e^{1,13\%.t}=2\Leftrightarrow 1,13\%.t=\ln_e 2\Leftrightarrow t=\dfrac{\ln 2}{1,13\%} \approx 61,34 \; (do\;  e > 1 \;).\;
		 \]
		Vậy sau $62$ năm tức đến năm $2085$ thì dân số ở địa phương đó sẽ gấp đôi dân số năm  $2023$.
	 }
\end{ex}

% nguồn: Ok_TraLoiNgan_Lop11_Lan1
\begin{ex}%[1D6V4-6]%[ma:749fdf47]
	Thực hiện một mẻ nuôi cấy vi khuẩn với $1\,000$ vi khuẩn ban đầu, nhà sinh học phát hiện ra số lượng vi khuẩn tăng thêm $25\%$ sau mỗi hai ngày. Công thức $P(t)=P_0\cdot a^t$ cho phép tính số lượng vi khuẩn của mẻ nuôi cấy sau $t$ ngày kể từ  thời điểm ban đầu. Sau bao nhiêu ngày thì số lượng vi khuẩn vượt gấp đôi số lượng ban đầu? Làm tròn kết quả đến hàng phần mười.
	 \shortans{6,2}
	 \loigiai{
	Theo  đề ta có $P(0)=1000$ và $P(2)=1000+25\%\cdot 1000=1250$.\\
	Từ  $P(0)=1000$ suy ra $P_0\cdot a^0=1000 \Rightarrow P_0=1000$.\\
	Từ $P(2)=1250$ suy ra $P_0\cdot  a^2=1250 \Rightarrow a^2=\dfrac{5}{4} \Rightarrow a= \dfrac{\sqrt{5}}{2}$.\\
	 Vậy $P(t)=1000\cdot \left(\dfrac{\sqrt{5}}{2}\right)^t$.\\
	Số  lượng vi khuẩn vượt gấp đôi số lượng ban đầu khi $P(t)=1000\cdot \left(\dfrac{\sqrt{5}}{2}\right)^t=2\cdot 1000 \Leftrightarrow \left(\dfrac{\sqrt{5}}{2}\right)^t=2 \Rightarrow t\approx 6{,}2$.\\
	Vậy sau $6{,}2$ ngày thì số lượng vi khuẩn vượt gấp đôi  số lượng ban đầu.
	 }
\end{ex}

% nguồn: Ok_TraLoiNgan_Lop11_Lan1
\begin{ex}%[1D6V4-6]%[ma:ab72dcc4]
	Ông Minh gửi ngân hàng số tiền $100$ triệu đồng với lãi suất $7\%$ một năm theo hình thức lãi kép. Biết rằng trong suốt quá trình gửi ông không rút tiền lãi. Hỏi sau ít nhất bao nhiêu  năm thì ông có nhiều hơn $500$ triệu?
	 \shortans{24}
	 \loigiai{
		Gọi $n$ là số năm ông Minh gửi, số tiền ông Minh khi gửi $n$ năm là $100 (1+0{,}07)^n$. Theo giả thiết ta  có
		\[100 (1+0{,}07)^n > 500 \Leftrightarrow n > \log_{1{,}07} 5 \approx  23{,}787\]
		Vậy ông Minh cần ít nhất $24$  năm.
	} 
\end{ex}

% nguồn: Ok_TraLoiNgan_Lop11_Lan2
\begin{ex}%[1D6V4-6]%[ma:30dee229]
	Áp suất không khí $P$ là một đại lượng được tính theo công thức $P=P_0 e^{-ix}$ trong đó $x$ là độ  cao, $P_0=760$ mmHg là áp suất ở mực nước biển, $i$ là hệ số suy giảm. Biết rằng, ở độ cao $1000$ $m$ thì áp suất của không khí là $672{,}72$ mmHg. Hỏi áp suất của không khí ở độ cao $15$ km là bao nhiêu?(Kết quả làm tròn đến hàng đơn vị)
	\shortans{122} 
	 \loigiai{Do ở độ cao $1000$ $m$, áp suất của không khí là $672{,}72$ mmHg nên ta có
		 $$
		672{,}72 =  760e^{-1000i}
		 \Leftrightarrow i = \dfrac{1}{1000} \ln \dfrac{672{,}72}{760}.
		 $$
		Khi ở độ cao $15$ km tức là $15000$ m thì  áp suất của không khí
		 $$P = 760e^{-15000 \times \dfrac{1}{1000} \ln \dfrac{672{,}72}{760}} \approx 122.$$
	Vậy, áp suất của không khí ở  độ cao $15$ km gần nhất với số $122$.}
\end{ex}

% nguồn: Ok_TraLoiNgan_Lop11_Lan2
\begin{ex}%[1D6V4-6]%[ma:b1958f3f]
	Ông An gửi $500$ triệu vào ngân hàng theo hình thức lãi kép trong một thời gian khá lâu với lãi suất ổn định trong suốt thời gian tiết kiệm là 10\% $1$  năm. Tết năm nay do dịch bệnh nên ông rút hết tiền trong ngân hàng ra để gia đình chi tiêu. Sau khi rút cả vốn lẫn lãi, ông trích ra $20$ triệu để sắm sửa đồ Tết thì ông còn $860$ triệu. Hỏi ông đã gửi tiết kiệm trong bao nhiêu năm?
	\shortans{6} 
	\loigiai{ 
		Giả  sử ông An đã gửi tiết kiệm trong $n$ năm. \\
		Số  tiền ông đã nhận được là $880$ triệu. \\
		Theo công thức lãi suất kép, ta  có
		 $$
		880 \cdot 10^6 = 500 \cdot 10^6 (1+0{,}1)^n 
		\Leftrightarrow n = \log_{1{,}1}  \dfrac{880}{500}
		\Leftrightarrow n  \approx 5{,}93.
		$$ 
		Vậy, ông A  đã gửi tiết kiệm trong $6$ năm.
	} 
\end{ex}

% nguồn: Ok_TraLoiNgan_Lop11_Lan2
\begin{ex}%[1D6V4-6]%[ma:3f0c6c11]
	Dân số thế giới  được ước tính theo công thức $P_n=P_0 e^{rn}$, trong đó $P_n$ là dân số của năm lấy làm mốc, $P_n$ là dân số sau $n$ năm, $r$ là tỉ lệ tăng dân số hằng năm. Biết rằng năm $2001$ dân số Việt Nam là $76$. $685$. $800$ người và tỉ lệ tăng dân số năm đó là $1$, 7\%. Hỏi cứ tăng dân số với tỉ lệ như vậy thì đến năm nào dân số nước ta ở mức $115$ triệu người
	 \shortans{2025}
	 \loigiai{
		Theo bài ra ta xét phương  trình
		 $$
		P_n =  115.10^6
		\Leftrightarrow P_0 e^{rn} =  115.10^6
		\Leftrightarrow  \ln P_0 + nr = \ln \left(115.10^6 \right).
		$$ 
		Suy ra $n = \dfrac{\ln \left(115.10^6 \right) -  \ln P_0}{r} \approx 23{,}8.$ \\
	Như vậy, đến năm $2025$ dân số nước ta sẽ ở mức $115$  triệu người.}
\end{ex}

% nguồn: Ok_TraLoiNgan_Lop11_Lan2
\begin{ex}%[1D6V4-6]%[ma:0b47e9eb]
	Trong Vật lí, sự phân rã của các chất phóng xạ được biểu diễn bởi công thức $m(t)=m_0 \left( \dfrac{1}{2} \right)^{\tfrac{t}{T}} $. Trong đó $m_0 $ là khối lượng ban đầu của chất phóng xạ, $T$ là chu kỳ bán rã. Một chất phóng xạ $X$ có chu kỳ bán rã là $T=20$ năm, người ta tìm được trong  một mẫu đồ cổ một lượng chất $X$ và xác định nó đã mất hơn $15\% $ lượng chất $X$ ban đầu của nó. Mẫu đồ cổ trên có tuổi ít nhất là bao nhiêu?
	 \shortans{5}
	\loigiai{ 
		Gọi khối lượng ban đầu của mẫu đồ cổ đó  là $m_0 $.\\
		Tại thời điểm $t$ tính từ thời điểm ban đầu ta có $m(t)=m_0 \left(\dfrac{1}{2}  \right)^{\tfrac{t}{T}} $. \\
		Theo giả thiết khối  lượng mẫu đồ cổ trên còn lại sẽ ít hơn $85\% $ so với ban đầu nghĩa là
		\begin{eqnarray*}
			m<85\% \cdot m_0 =\dfrac{17}{20} m_0 &\Rightarrow& m<\dfrac{17}{20} m_0\\ &\Rightarrow& m_0 \left(\dfrac{1}{2} \right)^{\tfrac{t}{T}} <\dfrac{17}{20} m_0\\  &\Leftrightarrow& \left(\dfrac{1}{2} \right)^{\tfrac{t}{T}} <\dfrac{17}{20}\\
			&\Rightarrow& \dfrac{t}{T}  >\log _{\tfrac{1}{2}} \dfrac{17}{20} \\
			&\Leftrightarrow& t>T\cdot  \log _{\tfrac{1}{2}} \dfrac{17}{20} \Rightarrow t>4{,}7\approx 5.
		\end{eqnarray*}
		Vậy  mẫu đồ cổ trên có tuổi ít nhất là $5$.
	 }
\end{ex}

% nguồn: Ok_TraLoiNgan_Lop11_Lan2
\begin{ex}%[1D6V4-6]%[ma:643a56e9]
	Ông B vay ngân hàng $600$ triệu đồng và trả góp hàng tháng. Cuối mỗi tháng bắt đầu từ tháng thứ nhất anh trả $10$ triệu đồng và chịu lãi suất $0$, 8\% trên tháng cho số tiền chưa trả. Với hình thức hoàn nợ  như vậy thì sao bao lâu ông B sẽ trả hết số nợ ngân hàng
	 \shortans{50}
	 \loigiai{
		Số tiền còn lại sau  $1$ tháng: $600 \cdot 10^6+600 \cdot 10^6 \cdot 0{,}8\%-10 \cdot 10^6$.\\
		Khi đó ta có thể gọi số tiền vay là $A$, lãi suất là $r$, số tiền trả mỗi cuối tháng là $m$ và $n$ là số tháng để  trả hết tiền. \\
		Vậy  số tiền còn nợ cuối tháng 1: $A+Ar-m=A(1+r) -m$ \\
		Số tiền còn nợ cuối  tháng 2: $A(1+r) -m+\left[A(1+r) -m \right]r-m=A(1+r)^2-\dfrac{m}{r}\left[(1+r)^2-1 \right]$ \\
		Số tiền còn nợ cuối tháng $n$:  $A(1+r)^n-\dfrac{m}{r}\left[(1+r)^n-1 \right]$ \\
		 Có nghĩa là khi trả hết tiền thì:
		\begin{align*}
			 A(1+r)^n-\dfrac{m}{r}\left[(1+r)^n-1 \right]=0
			&\Leftrightarrow 600  \cdot 10^6(1+0{,}8\%)^n-\dfrac{10 \cdot 10^6}{0{,}8}\left[(1+0{,}8\%)^n-1 \right]=0\\
			 &\Leftrightarrow(1+0{,}8\%)^n=1{,}48384 \\
			&\Rightarrow n=\log_{1+0{,}8\%}  1{,}48384 \approx 49{,}5
		\end{align*}
	Vậy, sau $50$ tháng thì ông  B trả hết nợ.}
\end{ex}

% nguồn: Ok_TraLoiNgan_Lop11_Lan2
\begin{ex}%[1D6V4-6]%[ma:5f68df36]
	Gọi $I(t)$ là số ca bị nhiễm bệnh Covid-19 ở quốc  gia $X$ sau $t$ ngày khảo sát. Khi đó ta có công thức $I(t)=A.e^{r_0(t-1) }$ với A là số ca bị nhiễm trong ngày khảo sát đầu tiên, $r_0$ là hệ số lây nhiễm. Biết rằng ngày đầu tiên khảo sát có $500$ ca bị nhiễm bệnh và ngày thứ $10$ khảo sát có $1000$ ca bị nhiễm bệnh. Hỏi ngày thứ $20$ số ca nhiễm bệnh gần nhất với số nào dưới đây, biết rằng trong suốt quá trình khảo sát hệ số lây nhiễm là không đổi?
	\shortans{2160} 
	\loigiai{ 
		Theo giả thiết ta có $I(1) = A = 500$. Ngày  thứ $10$ có $1000$ ca nên
		 $$
		 I(10) = A.e^{9r_0}
		\Leftrightarrow 1000  = 500.e^{9r_0}
		\Leftrightarrow r_0 = \ln \dfrac{2}{9}. 
		 $$
	Vậy, ngày thứ $20$, số  ca nhiễm bệnh là $I(20) = 500 \cdot e^{\tfrac{19 \ln 2}{9}} \approx 2160$.}
\end{ex}

% nguồn: Ok_TraLoiNgan_Lop11_Lan2
\begin{ex}%[1D6V4-6]%[ma:9926e87f]
	Anh Hải có $200\,000\,000$ đồng gửi ngân hàng kỳ hạn là $3$ tháng với lãi suất $0{,}55\% $ một tháng theo hình thức  lãi kỳ này được nhập vào vốn để tính lãi cho kỳ sau. Biết rằng nếu không rút lãi trong tất cả các quý định kỳ thì sau tối thiểu $n$ quý ($n\in \mathbb{N}^* $) gửi tiền vào ngân hàng anh Hải mới có số tiền lãi lớn hơn số tiền ban đầu gửi. Giá trị của $n$ là
	\shortans{43} 
	 \loigiai{
		 Ta có
		\begin{itemize}
			\item $r=3\cdot 0{,}55\% =1{,}65\% $ là lãi suất trong một kỳ  hạn;
			\item  $P_0 =200\,000\,000$ đồng là vốn ban đầu.
		\end{itemize}
		Tổng số tiền lãi và  vốn sau $n$ quý gửi tiền vào ngân hàng là $P_n =P_0 \cdot(1+1{,}65\%)^n $.\\
		Tổng số  tiền lãi sau $n$ quý gửi tiền vào ngân hàng là $P_n -P_0 =P_0 \cdot(1+1{,}65\%)^n -P_0 $.\\
		 Theo yêu cầu của bài toán ta có
		\begin{eqnarray*}
			P_n -P_0 >P_0 &\Leftrightarrow & P_0  \cdot (1+1{,}65\%)^n -P_0 >P_0\\
			&\Leftrightarrow&  P_0 \cdot(1+1{,}65\%)^n >2P_0 \\
			&\Leftrightarrow& n>\log  _{(1+1,65\%)} 2\approx 42{,}35.
		\end{eqnarray*}
	Do $n\in \mathbb{N}^* $ nên ta chọn $n=43$. Vậy sau  tối thiểu $43$ quý anh Hải mới có số tiền lãi lớn hơn số tiền ban đầu gửi}
\end{ex}

% nguồn: Ok_TraLoiNgan_Lop11_Lan2
\begin{ex}%[1D6V4-6]%[ma:cfd8b81c]
	Người ta nuôi cấy một loại vi khuẩn $V$ trong phòng thí nghiệm. Nếu số vi khuẩn ban đầu là $N_0$, tỉ lệ tăng trưởng vi khuẩn là $r$ thì sau $t$ giờ số vi khuẩn $N(t)$ nuôi cấy  được ước tính theo công thức $N(t)=N_0 e^{rt}$. Ban đầu có $300$ con vi khuẩn và sau $4$ giờ người ta thấy số lượng vi khuẩn đã tăng gấp đôi. Số giờ gần nhất để số lượng vi khuẩn thu được là $9\,000$ con là
	 \shortans{20}
	\loigiai{ 
		Áp dụng công thức $N(t) = N_0 e^{rt}$,  ta có
		$$ 
		600 =  300 \cdot e^{r \cdot 4}
		\Leftrightarrow 2 =  e^{4r}
		 \Leftrightarrow 4r = \ln 2
		\Leftrightarrow r =  \dfrac{\ln 2}{4}.
		 $$
		Để  số lượng vi khuẩn thu được là $9\,000$ con thì
		$$ 
		300 \cdot e^{\tfrac{\ln 2}{4} \cdot t} =  9\,000
		\Leftrightarrow e^{\tfrac{\ln 2}{4} \cdot  t} = 30
		\Leftrightarrow \dfrac{\ln 2}{4} \cdot t = \ln  30
		\Leftrightarrow t  \approx 19{,}628.
	 $$}
\end{ex}

% nguồn: Ok_TraLoiNgan_Lop11_Lan2
\begin{ex}%[1D6V4-6]%[ma:49f01886]
	Dân số thế giới được ước tính theo công thức $S=A e^{\ni}$, trong đó A là dân số của năm lấy mốc, $S$ là dân số sau $n$ năm, $i$ là tỷ lệ gia tăng dân số hàng năm. Biết năm $2023$ dân số thành phố Cần Thơ năm $2023$ ước tính là $1282000$ người và tỉ lệ gia tăng dân số là $1{,}03\%$. Hỏi đến năm bao nhiêu thì dân số thành  phố Cần Thơ đạt hơn $1{,}5$ triệu người?
	\shortans{} 
	\loigiai{Trả lời:$2039$ Lấy năm $2023$ làm mốc, khi đó $A = 1282000$. Giả sử sau $n$ năm thì dân số thành phố Cần Thơ đạt hơn $1500000$, tức là $1500000 = 1282000.e^{0{,}0103n} \Rightarrow e^{0{,}0103n} = \dfrac{750}{641} \Rightarrow 0{,}0103n = \ln \dfrac{750}{641}  \Rightarrow n \approx 15{,}25$ Có nghĩa là theo tốc độ tăng dân số này thì sau $16$ năm dân số thành phố sẽ đạt $1500000$ người vào năm $2039$.}
\end{ex}

% nguồn: Ok_TraLoiNgan_Lop11_Lan2
\begin{ex}%[1D6V4-6]%[ma:291acb7d]
	Bạn B hàng tháng gửi vào ngân hàng một số tiền như nhau là $2\,000\,000$ đồng, kỳ hạn là $1$ tháng với lãi suất $0{,}65\% $ một tháng. Sau tối thiểu $n$ tháng ($n\in \mathbb{N}^* $) số tiền cả vốn và lãi bạn B có được  lớn hơn $50\,000\,000$ đồng. Giá trị của $n$ là
	\shortans{24} 
	\loigiai{ 
		Gọi $a=2\,000\,000$ là số tiền hàng tháng bạn B gửi, $r$ là lãi  suất, ta có
		\begin{itemize}
			\item Số tiền bạn  B có ở cuối tháng thứ $1$: $P_1 =a+ar$;
			\item Số tiền bạn B có ở cuối tháng thứ $2$:  $P_2 =P_1 +a+(P_1 +a)r=a[(1+r)+(1+r)^2]$;
			\item Số tiền bạn  B có ở cuối tháng thứ $3$: $P_3 =P_2 +a+(P_2 +a)r=a[(1+r)+(1+r)^2 +(1+r)^3]$;
			\item Số tiền bạn B có ở cuối tháng thứ  $n$: $P_n =P_{n-1} +a+(P_{n-1} +a)r=a[(1+r)+(1+r)^2 +(1+r)^3 +\cdots+(1+r)^n]=\dfrac{a(1+r)[(1+r)^n -1]}{r} $.
		\end{itemize}
		Theo yêu câu  bài toán ta có
		\begin{eqnarray*}
			 \dfrac{a(1+r)[(1+r)^n -1]}{r} >25a &\Leftrightarrow& (1+r)^{n+1} >1+26r\\
			&\Leftrightarrow & n+1>\log _{1+r}  (1+26r)\\
			&\Leftrightarrow& n+1>\log  _{1+0{,}65\%} (1+26\cdot 0{,}65\%)\Leftrightarrow n>23{,}1.
		\end{eqnarray*}
		Do $n\in  \mathbb{N}^* \Rightarrow n=24$. \\
	Vậy sau tối thiểu $24$ tháng  bạn B sẽ có số tiền thỏa mãn yêu cầu bài toán.}
\end{ex}

% nguồn: Ok_TraLoiNgan_Lop11_Lan2
\begin{ex}%[1D6V4-6]%[ma:510ce691]
	Bạn An gửi vào ngân hàng số tiền $60\,000\,000$ đồng với kỳ hạn $1$ tháng và lãi suất $0{,}45\% $ một tháng theo hình thức lãi kỳ này được nhập vào vốn để tính lãi cho kỳ sau. Mỗi tháng An rút $1\,000\,000$ đồng vào ngày ngân hàng tính lãi để chi tiêu, sau tối thiểu $n$ tháng ($n\in \mathbb{N}^* $)  số tiền An còn lại không quá $40\,000\,000$ đồng. Giá trị của $n$ là
	 \shortans{26}
	\loigiai{ 
		Gọi $a=60\,000\,000$ là số tiền bạn An gửi, $r$ là lãi suất, $x$ đồng  là số tiền An rút mỗi tháng.
		\begin{itemize}
			\item Số tiền bạn An còn sau tháng thứ $1$: $P_1  =a+ar-x=a(r+1)-x$;
			\item Số tiền  bạn An còn sau tháng thứ $2$: $P_2 =P_1 +P_1 \cdot r-x=a(1+r)^2 -x(1+r)-x$;
			\item  Số tiền bạn An còn sau tháng thứ $3$: $P_3 =P_2 +P_2 \cdot r-x=a(1+r)^3 -x(1+r)^2 -x(1+r)-x$;
			\item Số tiền bạn An còn sau tháng thứ $n$: $P_n =P_{n-1}  +P_{n-1} \cdot r-x=a(1+r)^n -x(1+r)^{n-1} -\cdots-x(1+r)-x=(1+r)^n \left(a-\dfrac{x}{r}\right)+\dfrac{x}{r} $.
		\end{itemize}
		Theo yêu cầu bài toán  ta có
		\begin{eqnarray*}
			(1+r)^n \left(a-\dfrac{x}{r}\right)+\dfrac{x}{r} <40\,000\,000 &\Leftrightarrow& (1+r)^n \left(60-\dfrac{1}{0{,}45\%}\right)<40-\dfrac{1}{0{,}45\%} \\ 
			 &\Leftrightarrow& (1+r)^n >\dfrac{82}{73}\\
			&\Leftrightarrow& n>\log _{1+0{,}45\%} \dfrac{82}{73} \Leftrightarrow  n>25{,}89.
		\end{eqnarray*}
		Do $n\in \mathbb{N}^* \Rightarrow n=26$. Vậy sau tối thiểu $26$ tháng  bạn An sẽ còn lại số tiền không quá $40\,000\,000$ đồng.
	} 
\end{ex}

% nguồn: Ok_TraLoiNgan_Lop11_Lan2
\begin{ex}%[1D6V4-6]%[ma:4cebbead]
	Giả sử quá trình nuôi cấy vi khuẩn tuân theo quy luật tăng trưởng tự do. Khi đó, nếu gọi $N_0$ là số lượng vi khuẩn ban đầu và $N(t)$ là số lượng vi khuẩn sau $t$ giờ thì ta có $N(t)=N_0 e^{rt}$, trong đó $r$ là tỉ lệ tăng trưởng vi khuẩn mỗi giờ. Giả sử ban đầu có $500$ con vi khuẩn và sau $1$ giờ tăng  lên $800$ con. Tính số lượng vi khuẩn sau $3$ giờ.
	\shortans{2048} 
	 \loigiai{
		Ban đầu có $500$ con vi khuẩn và sau $1$  giờ tăng lên $800$ con nên ta có
		 $$
		800 =  500 \cdot e^r
		\Leftrightarrow e^r = 1{,}6 
		\Leftrightarrow r = \ln 1{,}6. 
		$$ 
		Số lượng vi  khuẩn tăng lên sau $3$ giờ là
		 $$
		N(3)  = 500 \cdot e^{3\ln 1{,}6} = 2048.
	 $$}
\end{ex}

% nguồn: Ok_TraLoiNgan_Lop11_Lan2
\begin{ex}%[1D6V4-6]%[ma:c74b2dd3]
	Sự tăng trưởng của một loại vi khuẩn trong phòng thí nghiệm được tính theo công thức $S(t)=S_0.e^{rt}$. Trong đó $S_0$ là số lượng vi khuẩn ban đầu,  $S(t)$ là số lượng vi khuẩn có sau $t$, $r$ là tỷ lệ tăng trưởng $(r > 0)$, $t$ là thời gian tăng trưởng. Biết rằng số lượng vi khuẩn ban đầu có $500$ con và sau $5$ giờ có $1\,500$ con. Hỏi sau bao nhiêu giờ kể từ lúc ban đầu có $500$ con để số lượng vi khuẩn đạt $121\,500$ con?
	\shortans{25} 
	 \loigiai{
		Ta có $S_0 = 500$; $5$ giờ =  $300$ phút. \\
		Sau $5$ giờ, số vi khuẩn là 
		$$ 
		 S(300) = 500 \cdot e^{300r}
		 \Leftrightarrow 1500 = 500.e^{300r}
		\Leftrightarrow  r = \dfrac{\ln 3}{300}.
		$$ 
		Vậy, khoảng thời gian $t$ kể từ lúc bắt đầu, có $500$ con vi khuẩn  đến khi số lượng vi khuẩn đạt $121\,500$ con thỏa mãn
		$$ 
		 121500 = 500 \cdot e^{rt}
		\Leftrightarrow t = \dfrac{\ln 243}{r}  = \dfrac{300 \ln 243}{\ln 3} = 1500 = 25.
	$$} 
\end{ex}

% nguồn: Ok_TraLoiNgan_Lop11_Lan2
\begin{ex}%[1D6V4-6]%[ma:d2f12f0e]
	Cho áp suất không khí $P$ suy giảm mũ so với độ cao $x$, tức $P$ giảm theo công thức $P=P_0 e^{-ix}$ trong đó $P_0=760 \text{mmHg}$ là áp suất ở mực nước biển $(x=0)$, $i$ là hệ số suy giảm. Biết rằng ở độ cao $1000 m$ thì áp suất của không khí là $672{,}71 \text{mmHg}$. Hỏi áp suất không khí ở độ cao $3580 m$ là bao nhiêu?(Kết quả  làm tròn đến hàng đơn vị)
	\shortans{491} 
	 \loigiai{
		Áp dụng công thức  $P=P_0 e^{-ix}$, ở độ cao $1000$ m, ta có
		$P_0=760$ mmHg, $x=1000$ m, $P=672{,}71$ mmHg, từ giả thiết này ta tìm được hệ số suy  giảm $i$. \\
		Ta  có
		$$ 
		672{,}71 = 760e^{-1000i} 
		 \Leftrightarrow 1000i = \ln \dfrac{672{,}71}{760}
		\Leftrightarrow i = \dfrac{1}{1000} \ln \dfrac{672{,}71}{760}. 
		 $$
		Khi đó ở  độ cao $3580$ m, áp suất của không khí là
		 $$
		P = 760e^{-\dfrac{1}{1000} \ln \dfrac{672{,}71}{760} \times  3580} \approx 491{,}04.
		 $$
	 }
\end{ex}

% nguồn: Ok_TraLoiNgan_Lop11_Lan1
\begin{ex}%[1D6C4-6]%[ma:7ea880d2]
	Sau khi tốt nghiệp đại học, anh Nam thực hiện một dự án khởi nghiệp. Anh vay vốn từ ngân hàng $200$ triệu đồng với lãi suất $0,6\%$ một tháng. Phương án trả nợ của Nam là sau đúng một tháng kể từ thời điểm vay anh tao bắt đầu trả nợ, hai lần liên tiếp cách nhau đúng một tháng, số tiền trả của mỗi lần là như nhau và hoàn thành sau đúng $5$ năm kể từ khi vay. Tuy nhiên sau khi dự án có hiệu quả và trả nợ được $12$ tháng theo phương án cũ, anh Nam muốn rút ngắn thời gian trả nợ nên từ tháng tiếp theo, mỗi tháng anh trả nợ cho ngân hàng $9$ triệu đồng. Biết rằng ngân hàng chỉ tính tiền lãi trên số dư  nợ thực tế của tháng đó. Hỏi sau ít nhất bao nhiêu tháng từ thời điểm vay anh Nam trả hết nợ?
	\shortans{32} 
	 \loigiai{
		 Gọi $m$ là số tiền dự định trả hằng tháng.\\
		$A$ là  số tiền gốc đã vay.\\
		$i$ là lãi suất 1 tháng.\\ 
		 Theo dự định số tiền phải trả hàng tháng bằng
		 $$m=200\cdot 0{,}006\cdot \dfrac{\left( 1,006 \right)^{60}}{\left( 1,006 \right)^{60}-1} =\dfrac{6}{5}\cdot \dfrac{\left( 1,006 \right)^{60}}{\left( 1,006 \right)^{60}-1}.$$
		Số  tiền còn nợ sau 12 tháng là\\
		\centerline 
		{$\begin{aligned}
		A_{12}=&\ A\cdot \left( 1+i \right)^{12}-\left[m+m\left( 1+i \right)+\cdots+m\left( 1+i \right)^{11} \right]\\ 
		=&\ A\left( 1+i \right)^{12}-m\cdot \dfrac{\left( 1+i \right)^{12}-1}{i}=165{,}53  \text{ (triệu đồng)}.
	\end{aligned}$}\\
	Gọi $n$ là  số tháng tính từ thời điểm hết 12 tháng đến lúc trả hết nợ có\\
	\centerline{ 
	$\begin{aligned}
	&\ A_{12}\cdot (1+i)^n -\left[9+9\cdot(1+i) +\cdots +9(1+i)^{n-1} \right]=0 \\ 
	\Leftrightarrow &\ A_{12}\cdot (1+i)^n -9\cdot  \dfrac{(1+i)^n-1}{i}\le 0 \\
	\Leftrightarrow &\ (1{,}006)^n  \left( 1500-A_{12} \right)\le 1500\\
	\Leftrightarrow &\ n\ge  \log_{1,006}\dfrac{1500}{1500-A_{12}} \approx 19{,}55 .
	\end{aligned}$	}\\
Vậy  anh Nam cần ít nhất $12+20=32$ tháng để trả hết nợ.}
\end{ex}


%%%%%%%%%%%%%%%% TỰ LUẬN %%%%%%%%%%%%%%%%

% nguồn: Ok_TuanLuan_Toan11_Lan1
\begin{bt}%[1D6H4-6]%[5-KNTT]%[ma:33bd75a2]
	Số lượng của một loài vi khuẩn trong phòng thí nghiệm được tính theo công thức $S(t)=A\cdot e^{rt}$, trong đó $A$ là số lượng vi khuẩn ban đầu, $S(t)$ là số lượng vi khuẩn có sau $t$ (phút), $r$  là tỉ lệ tăng trưởng $(r>0), t$ (tính theo phút) là thời gian tăng trưởng. Biết rằng số lượng vi khuẩn ban đầu có $500$ con và sau $6$ giờ có $2000$ con. Hỏi ít nhất bao nhiêu giờ, kể từ lúc bắt đầu, số lượng vi khuẩn đạt ít nhất $120000$ con?
	 \loigiai{
		Ta có $A=500$,  $S(360)=2000$, $6$ giờ $=360$ phút.\\
		Sau $6$ giờ số lượng vi  khuẩn là $2000$ con, tức là $2000=500\cdot e^{r\cdot 360}$.\\
		Số lượng vi khuẩn  đạt ít nhất $120000$ con, nghĩa là $500\cdot e^{\frac{\ln 4}{360}\cdot t} \geq 120000$ $\Leftrightarrow e^{\frac{\ln 4}{360}\cdot t} \geq 240 \Leftrightarrow \frac{\ln 4}{360}\cdot t \geq \ln 240 \Leftrightarrow t \geq \frac{360\cdot \ln 240}{\ln 4} \approx 1423,24$ (phút).\\
		Vậy sau ít nhất $24$ (giờ) thì số lượng  vi khuẩn đạt ít nhất $120000$ con.
	 }
\end{bt}

% nguồn: Ok_TuanLuan_Toan11_Lan1
\begin{bt}%[1D6H4-6]%[ma:9b844a2d]
	Trong nông nghiệp bèo hoa dâu được dùng làm phân bón, nó rất tốt cho cây trồng. Mới đây, các nhà khoa học Việt Nam đã phát hiện ra  bèo hoa dâu có thể dùng để chiết xuất ra chất có tác dụng kích thích hệ miễn dịch và hỗ trợ điều trị bệnh ung thư. Bèo hoa dâu được thả nuôi trên mặt nước. Một người đã thả một lượng bèo hoa dâu chiếm $4\%$ diện tích mặt hồ. Biết rằng cứ sau đúng một tuần bèo phát triển thành 3 lần số lượng đã có và giả sử tốc độ phát triển của bèo ở mọi thời điểm như nhau. Hỏi sau ít nhất bao nhiêu ngày bèo sẽ vừa phủ kín mặt hồ?
	 \loigiai{
		 Số lượng bèo ban đầu chiếm $0{,}04$ diện tích mặt hồ.
		\begin{itemize}
			\item Sau $1$ tuần số lượng bèo là $0{,}04\times  3$ diện tích mặt hồ.
			\item Sau $2$ tuần số lượng bèo là $0{,}04\times 3^2$ diện tích mặt hồ. 
			\item Sau $n$ tuần số  lượng bèo là $0{,}04\times 3^n$ diện tích mặt hồ.
		\end{itemize}
		Để bèo phủ kín mặt hồ thì $0{,}04\times 3^n=1\Rightarrow 3^n=25\Rightarrow n=\log_325$ (tuần).\\ 
		Số ngày tương ứng là $7n=7\log_325\approx 20{,}51$  (ngày).\\
		Vậy sau ít nhất $21$ ngày thì bèo  hoa dâu sẽ phủ kín mặt hồ.
	 }
\end{bt}

% nguồn: Ok_TuanLuan_Toan11_Lan1
\begin{bt}%[1D6H4-6]%[ma:637c8a9d]
	Trong một  trận động đất, năng lượng giải toả $E$ (đơn vị: Jun, kí hiệu J) tại tâm địa chấn ở $M$ độ Richter được xác định xấp xỉ bởi công thức
	$$ 
	\log E \approx 11{,}4  + 1{,}5 \cdot M
	$$ 
	Tính xấp xỉ năng lượng giải toả tại tâm  địa chấn ở $3$ độ Richter.
	 \loigiai{
		Thay $M = 3$ vào công thức $\log E \approx 11{,}4 + 1{,}5 \cdot M$, ta có năng lượng giải tỏa tại tâm địa chấn ở  $3$ độ Richter là
		$$\log E \approx 11{,}4  + 1{,}5 \cdot 3 = 15{,}9 \Leftrightarrow E = 10^{15{,9}}.$$
		Vậy năng lượng giải toả tại  tâm địa chấn ở $3$ độ Richter là $E = 10^{15{,9}} \approx 7{,}94 \cdot 10^{15}$ J.
	} 
\end{bt}

% nguồn: Ok_TuanLuan_Toan11_Lan1
\begin{bt}%[1D6H4-6]%[ma:86d4a4a7]
	Độ $pH$ của một dung dịch được tính theo  công thức $pH=-\log\left[H^+ \right]$, trong đó $\left[H^+ \right]$ là nồng độ $H^+$ của dung dịch đó tính bằng mol/lít. Biết dung dịch A có độ $pH$ bằng $2{,}1$. Hỏi nồng độ $H^+$ trong dung dịch A là bao nhiêu?
	\loigiai{ 
	Ta có $$-\log\left[H^+  \right]=2{,}1\Rightarrow \left[H^+ \right]=10^{-2{,}1}.$$
	 }
\end{bt}

% nguồn: Ok_TuanLuan_Toan11_Lan1
\begin{bt}%[1D6H4-6]%[ma:3bbabf14]
	[1,0 điểm] Người ta có thể gửi tiền vào ngân hàng với thể thức lãi suất kép theo định kì. Theo thể thức này, nếu đến cuối kì hạn người gửi không  rút tiền lãi ra thì tiền lãi được tính vào vốn để tính lãi suất cho kì hạn kế tiếp.
	 Nếu một người gửi số tiền $A$ với lãi suất $r$ cho mỗi kì hạn thì sau $n$ kì hạn, số tiền người ấy thu được cả vốn lẫn lãi là $T=A(1+r)^n$.
	Năm $2024$, Ông An gửi tiền vào một ngân hàng theo hình thức lãi suất kép với kì hạn $1$ năm có lãi suất $5\%/$ năm. Hỏi đến năm bao nhiêu thì Ông An thu được cả vốn và lãi gấp đôi số tiền gửi ban đầu. Biết rằng lãi suất không thay đổi qua các năm và Ông An không rút tiền ra trong  suốt quá trình gửi.
	\loigiai{ 
		Áp dụng công thức lãi kép với lãi suất $r=5\%/$ năm thì số tiền cả vốn và lãi sau  $n$ năm là $T=A(1+0{,}05)^n$.\\
		Số tiền ông  An thu được gấp đôi số tiền ban đầu khi $$A(1+0{,}05)^n=2A \Rightarrow n=\log_{1{,}05}2 \approx 14{,}2.$$
		Vậy sau  $15$ năm, nghĩa là đến năm $2039$ thì ông An thu được cả vốn và lãi gấp đôi số tiền gửi ban đầu.
	} 
\end{bt}

% nguồn: Ok_TuanLuan_Toan11_Lan1
\begin{bt}%[1D6H4-6]%[ma:c680e744]
	Một người gửi ngân hàng 500 triệu đồng với kì hạn 1 tháng theo hình thức lãi kép, lãi suất  $0,6\%$ một tháng (kể từ tháng thứ hai trở đi, tiền lãi được tính theo phần trăm của tồng tiền gốc và tiền lãi tháng trước đó). Hỏi sau ít nhất bao nhiêu tháng thì người đó có tối thiểu 600 triệu đồng trong tài khoản tiết kiệm, biết rằng ngân hàng chi tính lãi khi đến kì hạn?
	\loigiai{ 
	Theo hình thức lãi kép, tổng số tiền cà gốc lẫn lãi trong tài khoàn của người đó sau $n$ tháng là: $T=A(1+r)^n=500(1+0,6\%)^n=500\cdot 1,006^n$  (triệu đồng).\\
	Theo đề bài $T\ge  600\Rightarrow 500\cdot 1,006^n\ge 600\Leftrightarrow 1,006^n\ge\dfrac{6}{5}\Leftrightarrow n\ge\log_{1,006}\dfrac{6}{5}\approx 30,5$.\\
	Vì ngân hàng chi tính lãi khi đến kì hạn nên phải sau  it nhất 31 tháng người đó mới có tối thiểu 600 triệu đồng trong tài khoản.\\
	} 
\end{bt}

% nguồn: Ok_TuanLuan_Toan11_Lan1
\begin{bt}%[1D6H4-6]%[5-KNTT]%[ma:9e7431e1]
	Số lượng vi khuẩn $V$ trong phòng thí nghiệm tính theo công thức $s(t)=s_0\cdot 2^t$ trong đó $s_0$ là số lượng vi khuẩn $V$ lúc đầu, $s(t)$ là số lượng vi khuẩn có trong $t$ phút. Biết sau 3 phút thì số lượng  vi khuẩn $A$ là $625$ nghìn con. Hỏi sau $9$ phút thì số lượng vi khuẩn $V$ bao nhiêu?
	\loigiai{ 
		Vì sau $3$ phút thì số lượng vi khuẩn  $A$ là $625$ nghìn con nên $625000=s_0\cdot 2^3$.\\
		Số lượng vi khuẩn $V$ sau $9$ phút là $s(t)=\dfrac{625000}{2^3}\cdot 2^9=625000\cdot 2^6=4\cdot 10^7$  (con).
	} 
\end{bt}

% nguồn: Ok_TuanLuan_Toan11_Lan1
\begin{bt}%[1D6H4-6]%[ma:a7c876c9]
	Giả sử sự tăng trưởng của một loại vi khuẩn trong quá trình nuôi cấy tuân theo công thức $ N(t) = N_{0}\cdot e^{rt} $ , trong đó $ N_{0} $ là số lượng vi khuẩn ban đầu,$ r $ là tỉ lệ tăng trưởng $ (r > 0) $, $ t $ là thời gian tăng trưởng. Biết rằng số lượng vi khuẩn ban đầu là $ 500 $ con, và sau $ 2 $ giờ có $ 1500 $con. Hỏi sau bao lâu thì số lượng vi khuẩn ban đầu sẽ tăng  gấp đôi?
	 \loigiai{
		Tỉ lệ tăng trưởng của loài vi khuẩn này là $ 1500 =  500 \cdot e^{2r} \Leftrightarrow e^{2r} = 3 \Leftrightarrow 2r = \ln 3 \Leftrightarrow r = \dfrac{\ln 3}{2} $. \\
		Thời gian để số lượng vi khuẩn ban đầu tăng gấp đôi là \begin{center}
		$ 2N_{0} =  N_{0} \cdot e^{rt} \Rightarrow e^{rt} = 2 \Rightarrow rt = \ln 2 \Rightarrow t = \dfrac{2\ln 2}{ln3} \approx 1{,}26 $ (giờ).
	\end{center}
} 
\end{bt}

% nguồn: Ok_TuanLuan_Toan11_Lan1
\begin{bt}%[1D6H4-6]%[ma:c449e64b]
	Đồng vị phóng xạ Uranium-235  (thường được sử dụng trong điện hạt nhân) có chu kì bán rã là $T=703\,800\,000$ năm. Theo đó, nếu ban đầu có $100$ gam Uranium-235 thì sau $t$ năm, do bị phân rã, lượng Uranium-235 còn lại được tính theo công thức $M=100\cdot\left(\dfrac{1}{2} \right)^{\tfrac{t}{T}}$ (gam). Sau thời gian bao lâu thì lượng Uranium-235 còn lại không quá $90\%$ so với ban đầu.
	 \loigiai{
	Theo đề bài ta có $$\left(\dfrac{1}{2} \right)^{\tfrac{t}{T}}\le 0{,}9\Leftrightarrow\dfrac{t}{T}\ge\log_{\tfrac{1}{2}}0{,}9\Leftrightarrow  t\ge T\cdot \log_{\tfrac{1}{2}}0{,}9\approx 106\,979\,777\text{ (năm)}.$$
	Vậy sau ít nhất $106\,979\,777\text{ (năm)}$ thì thỏa yêu cầu bài toán. 
	 }
\end{bt}

% nguồn: Ok_TuanLuan_Toan11_Lan1
\begin{bt}%[1D6H4-6]%[ma:f728a27c]
	Ông Quý dự định  gửi vào ngân hàng một số tiền với lãi suất $6{,}5\%$ một năm. Biết rằng, cứ sau mỗi năm số tiền lãi sẽ được nhập vào vốn ban đầu. Tính số tiền tối thiểu $x$ (triệu đồng, $x\in \mathbb{N}$) ông Quý gửi vào ngân hàng để sau $3$ năm số tiền lãi đủ mua một chiếc xe gắn máy trị giá 30 triệu đồng.
	 \loigiai{
		Công thức lãi kép $T=A(1+r)^n$.\\ 
		Tiền lãi ông Quý có  sau $3$ năm sẽ là tiền gốc lẫn lãi trừ đi số tiền gốc ban đầu.\\
		Ta có $A(1+6{,}5\%)^3-A\ge 30\Leftrightarrow A\ge \dfrac{30}{(1+6{,}5\%)^3-1} \approx  144{,}26$ triệu.
	 }
\end{bt}

% nguồn: Ok_TuanLuan_Toan11_Lan1
\begin{bt}%[1D6H4-6]%[ma:4cab3427]
	Với nước biển có nồng độ muối $30\%$, nhiệt độ  $T$ ($^\circ$C) của nước biển được tính bởi công thức $T = 7{,}9\ln (1{,}0245 - d) + 61{,}84$ trong đó $d$ (g/cm$^3$) là khối lượng riêng của nước biển.
	\begin{flushright}
		(Ron  Larson, Intermediate Algebra, Cengage)
	\end{flushright}
	Biết vùng biển khơi mặt ở một khu vực có nồng độ muối $30\%$ và  nhiệt độ là $8^\circ$ C. Tính khối lượng riêng của nước biển ở vùng biển đó (làm tròn đến $4$ chữ số thập phân sau dấu phẩy).
	 \loigiai{
		Vì vùng biển khơi mặt ở một khu vực có nồng độ muối $30\%$ và nhiệt độ là  $8^\circ$ C nên $T = 8$.\\
		Thay $T  = 8$ vào công thức $T = 7{,}9\ln (1{,}0245 - d) + 61{,}84$, ta được
		\begin{eqnarray*}
			& &T = 7{,}9\ln (1{,}0245  - d) + 61{,}84\\
			&\Leftrightarrow& 8 = 7{,}9\ln (1{,}0245 - d) + 61{,}84\\ 
			&\Leftrightarrow& 7{,}9\ln (1{,}0245 -  d) = -53{,}84\\
			&\Leftrightarrow& \ln  (1{,}0245 - d) = -\dfrac{2692}{395}\\
			 &\Leftrightarrow& 1{,}0245 - d = \mathrm{e}^{-\frac{2692}{395}}\\
			&\Leftrightarrow& d  = 1{,}0245 - \mathrm{e}^{-\frac{2692}{395}}\\
			&\Leftrightarrow&  d \approx 1{,}0234.
		\end{eqnarray*}
		Vậy khối lượng riêng của nước biển ở vùng biển đó  là $1{,}0234$ g/cm$^3$.
	 }
\end{bt}

% nguồn: Ok_TuanLuan_Toan11_Lan2
\begin{ex}%[1D6H4-6]%[5-KNTT]%[ma:5e1eb11f]
	Số lượng của một loài vi khuẩn trong phòng thí nghiệm được tính theo công thức $S(t)=A\cdot e^{rt}$, trong đó $A$ là số lượng vi khuẩn ban đầu, $S(t)$ là số lượng vi khuẩn có  sau $t$ (phút), $r$ là tỉ lệ tăng trưởng $(r>0), t$ (tính theo phút) là thời gian tăng trưởng. Biết rằng số lượng vi khuẩn ban đầu có $500$ con và sau $6$ giờ có $2000$ con. Hỏi ít nhất bao nhiêu giờ, kể từ lúc bắt đầu, số lượng vi khuẩn đạt ít nhất $120000$ con?
	 \loigiai{
		Ta có $A=500$, $S(360)=2000$, $6$ giờ  $=360$ phút.\\
		Sau $6$ giờ số lượng vi khuẩn là  $2000$ con, tức là $2000=500\cdot e^{r\cdot 360}$.\\
		Số lượng vi khuẩn đạt ít nhất $120000$ con, nghĩa là $500\cdot e^{\frac{\ln 4}{360}\cdot t} \geq 120000$ $\Leftrightarrow e^{\frac{\ln 4}{360}\cdot t} \geq  240 \Leftrightarrow \frac{\ln 4}{360}\cdot t \geq \ln 240 \Leftrightarrow t \geq \frac{360\cdot \ln 240}{\ln 4} \approx 1423,24$ (phút).\\
		Vậy sau ít nhất $24$ (giờ) thì số lượng vi khuẩn  đạt ít nhất $120000$ con.
	 }
\end{ex}

% nguồn: Ok_TuanLuan_Toan11_Lan2
\begin{ex}%[1D6H4-6]%[ma:f77ca23b]
	Trong nông nghiệp bèo hoa dâu được dùng làm phân bón, nó rất tốt cho cây trồng. Mới đây, các nhà khoa học Việt Nam đã phát hiện ra bèo hoa dâu có thể  dùng để chiết xuất ra chất có tác dụng kích thích hệ miễn dịch và hỗ trợ điều trị bệnh ung thư. Bèo hoa dâu được thả nuôi trên mặt nước. Một người đã thả một lượng bèo hoa dâu chiếm $4\%$ diện tích mặt hồ. Biết rằng cứ sau đúng một tuần bèo phát triển thành 3 lần số lượng đã có và giả sử tốc độ phát triển của bèo ở mọi thời điểm như nhau. Hỏi sau ít nhất bao nhiêu ngày bèo sẽ vừa phủ kín mặt hồ?
	 \loigiai{
		 Số lượng bèo ban đầu chiếm $0{,}04$ diện tích mặt hồ.
		\begin{itemize}
			\item Sau $1$ tuần số lượng bèo là $0{,}04\times 3$ diện tích mặt hồ. 
			\item Sau $2$ tuần số lượng bèo là $0{,}04\times 3^2$ diện tích mặt  hồ.
			\item Sau $n$ tuần  số lượng bèo là $0{,}04\times 3^n$ diện tích mặt hồ.
		\end{itemize}
		Để bèo  phủ kín mặt hồ thì $0{,}04\times 3^n=1\Rightarrow 3^n=25\Rightarrow n=\log_325$ (tuần).\\
		Số ngày tương ứng là  $7n=7\log_325\approx 20{,}51$ (ngày).\\
		Vậy sau ít nhất  $21$ ngày thì bèo hoa dâu sẽ phủ kín mặt hồ.
	 }
\end{ex}

% nguồn: Ok_TuanLuan_Toan11_Lan2
\begin{ex}%[1D6H4-6]%[ma:0f281516]
	Trong một trận động đất, năng lượng giải toả $E$ (đơn vị: Jun, kí hiệu J) tại tâm địa chấn ở $M$ độ Richter được xác định xấp xỉ bởi công thức 
	 $$
	\log  E \approx 11{,}4 + 1{,}5 \cdot M
	 $$
	Tính xấp xỉ năng lượng giải toả tại tâm địa chấn ở $3$ độ Richter. 
	\loigiai{ 
		Thay $M = 3$ vào công thức $\log E \approx 11{,}4 + 1{,}5 \cdot M$, ta  có năng lượng giải tỏa tại tâm địa chấn ở $3$ độ Richter là
		$$\log E \approx 11{,}4 + 1{,}5 \cdot 3 = 15{,}9  \Leftrightarrow E = 10^{15{,9}}.$$
		Vậy năng lượng giải toả tại tâm địa chấn ở $3$ độ Richter là $E  = 10^{15{,9}} \approx 7{,}94 \cdot 10^{15}$ J.
	} 
\end{ex}

% nguồn: Ok_TuanLuan_Toan11_Lan2
\begin{ex}%[1D6H4-6]%[ma:a6de01be]
	Độ $pH$ của một dung dịch được tính theo công thức $pH=-\log\left[H^+ \right]$, trong đó $\left[H^+ \right]$ là nồng độ $H^+$ của dung dịch đó tính bằng mol/lít. Biết dung dịch A có độ $pH$ bằng  $2{,}1$. Hỏi nồng độ $H^+$ trong dung dịch A là bao nhiêu?
	 \loigiai{
	Ta có $$-\log\left[H^+ \right]=2{,}1\Rightarrow \left[H^+ \right]=10^{-2{,}1}.$$ 
	 }
\end{ex}

% nguồn: Ok_TuanLuan_Toan11_Lan2
\begin{ex}%[1D6H4-6]%[ma:50110553]
	[1,0 điểm] Người ta có thể gửi tiền vào ngân hàng với thể thức lãi suất  kép theo định kì. Theo thể thức này, nếu đến cuối kì hạn người gửi không rút tiền lãi ra thì tiền lãi được tính vào vốn để tính lãi suất cho kì hạn kế tiếp.
	Nếu một người gửi số tiền $A$ với lãi suất $r$ cho mỗi kì hạn thì sau $n$ kì hạn, số tiền người ấy thu được cả vốn lẫn lãi là $T=A(1+r)^n$. 
	Năm $2024$, Ông An gửi tiền vào một ngân hàng theo hình thức lãi suất kép với kì hạn $1$ năm có lãi suất $5\%/$ năm. Hỏi đến năm bao nhiêu thì Ông An thu được cả vốn và lãi gấp đôi  số tiền gửi ban đầu. Biết rằng lãi suất không thay đổi qua các năm và Ông An không rút tiền ra trong suốt quá trình gửi.
	\loigiai{ 
		Áp dụng công thức lãi kép với lãi suất $r=5\%/$ năm  thì số tiền cả vốn và lãi sau $n$ năm là $T=A(1+0{,}05)^n$.\\
		Số  tiền ông An thu được gấp đôi số tiền ban đầu khi $$A(1+0{,}05)^n=2A \Rightarrow n=\log_{1{,}05}2 \approx 14{,}2.$$
		Vậy sau $15$ năm, nghĩa là đến năm $2039$  thì ông An thu được cả vốn và lãi gấp đôi số tiền gửi ban đầu.
	} 
\end{ex}

% nguồn: Ok_TuanLuan_Toan11_Lan2
\begin{ex}%[1D6H4-6]%[ma:d42f236b]
	Một người gửi ngân hàng 500 triệu đồng với kì hạn 1 tháng theo  hình thức lãi kép, lãi suất $0,6\%$ một tháng (kể từ tháng thứ hai trở đi, tiền lãi được tính theo phần trăm của tồng tiền gốc và tiền lãi tháng trước đó). Hỏi sau ít nhất bao nhiêu tháng thì người đó có tối thiểu 600 triệu đồng trong tài khoản tiết kiệm, biết rằng ngân hàng chi tính lãi khi đến kì hạn?
	\loigiai{ 
	Theo hình thức lãi kép, tổng số tiền cà gốc lẫn lãi trong tài khoàn  của người đó sau $n$ tháng là: $T=A(1+r)^n=500(1+0,6\%)^n=500\cdot 1,006^n$ (triệu đồng).\\
	Theo đề bài $T\ge 600\Rightarrow  500\cdot 1,006^n\ge 600\Leftrightarrow 1,006^n\ge\dfrac{6}{5}\Leftrightarrow n\ge\log_{1,006}\dfrac{6}{5}\approx 30,5$.\\
	Vì ngân hàng chi tính lãi khi đến kì hạn nên  phải sau it nhất 31 tháng người đó mới có tối thiểu 600 triệu đồng trong tài khoản.\\
	 }
\end{ex}

% nguồn: Ok_TuanLuan_Toan11_Lan2
\begin{ex}%[1D6H4-6]%[5-KNTT]%[ma:69261cbf]
	Số lượng vi khuẩn $V$ trong phòng thí nghiệm tính theo công thức $s(t)=s_0\cdot 2^t$  trong đó $s_0$ là số lượng vi khuẩn $V$ lúc đầu, $s(t)$ là số lượng vi khuẩn có trong $t$ phút. Biết sau 3 phút thì số lượng vi khuẩn $A$ là $625$ nghìn con. Hỏi sau $9$ phút thì số lượng vi khuẩn $V$ bao nhiêu?
	 \loigiai{
		Vì sau $3$ phút thì số lượng vi khuẩn $A$ là $625$ nghìn con nên $625000=s_0\cdot 2^3$.\\ 
		Số lượng vi khuẩn $V$ sau $9$ phút là $s(t)=\dfrac{625000}{2^3}\cdot 2^9=625000\cdot 2^6=4\cdot 10^7$  (con).
	 }
\end{ex}

% nguồn: Ok_TuanLuan_Toan11_Lan2
\begin{ex}%[1D6H4-6]%[ma:1264c61c]
	Giả sử sự tăng trưởng của một loại vi khuẩn trong quá trình nuôi cấy tuân theo công thức $ N(t) = N_{0}\cdot  e^{rt} $ , trong đó $ N_{0} $ là số lượng vi khuẩn ban đầu,$ r $ là tỉ lệ tăng trưởng $ (r > 0) $, $ t $ là thời gian tăng trưởng. Biết rằng số lượng vi khuẩn ban đầu là $ 500 $ con, và sau $ 2 $ giờ có $ 1500 $con. Hỏi sau bao lâu thì số lượng vi khuẩn ban đầu sẽ tăng gấp đôi?
	 \loigiai{
		Tỉ lệ tăng trưởng của loài vi khuẩn này là $ 1500  = 500 \cdot e^{2r} \Leftrightarrow e^{2r} = 3 \Leftrightarrow 2r = \ln 3 \Leftrightarrow r = \dfrac{\ln 3}{2} $. \\
		Thời gian để số lượng vi khuẩn ban đầu tăng gấp đôi là \begin{center}
		$ 2N_{0} = N_{0} \cdot e^{rt} \Rightarrow e^{rt} = 2 \Rightarrow rt = \ln 2 \Rightarrow t = \dfrac{2\ln 2}{ln3} \approx 1{,}26 $  (giờ).
	\end{center}
 }
\end{ex}

% nguồn: Ok_TuanLuan_Toan11_Lan2
\begin{ex}%[1D6H4-6]%[ma:5eec081e]
	Đồng vị phóng xạ Uranium-235 (thường được sử dụng trong điện hạt nhân) có chu kì bán rã là $T=703\,800\,000$ năm. Theo đó, nếu ban đầu có $100$ gam Uranium-235 thì sau $t$ năm, do bị phân rã, lượng Uranium-235 còn lại được tính theo công thức $M=100\cdot\left(\dfrac{1}{2} \right)^{\tfrac{t}{T}}$ (gam). Sau thời  gian bao lâu thì lượng Uranium-235 còn lại không quá $90\%$ so với ban đầu.
	\loigiai{ 
	Theo đề bài ta có  $$\left(\dfrac{1}{2} \right)^{\tfrac{t}{T}}\le 0{,}9\Leftrightarrow\dfrac{t}{T}\ge\log_{\tfrac{1}{2}}0{,}9\Leftrightarrow t\ge T\cdot \log_{\tfrac{1}{2}}0{,}9\approx 106\,979\,777\text{ (năm)}.$$
	Vậy sau ít nhất $106\,979\,777\text{ (năm)}$ thì thỏa yêu cầu bài toán. 
	} 
\end{ex}

% nguồn: Ok_TuanLuan_Toan11_Lan2
\begin{ex}%[1D6H4-6]%[ma:4dbf34f1]
	Ông Quý dự định gửi vào ngân hàng một số tiền với lãi suất $6{,}5\%$ một năm. Biết rằng, cứ sau mỗi năm số tiền lãi sẽ được nhập vào vốn ban đầu. Tính số tiền tối thiểu $x$ (triệu đồng, $x\in \mathbb{N}$) ông Quý gửi vào ngân hàng để sau $3$ năm số tiền lãi đủ mua một chiếc xe gắn máy trị giá 30 triệu đồng. 
	 \loigiai{
		Công thức  lãi kép $T=A(1+r)^n$.\\
		Tiền lãi ông Quý có sau $3$ năm sẽ là tiền gốc lẫn lãi trừ đi số tiền  gốc ban đầu.\\
		Ta có $A(1+6{,}5\%)^3-A\ge 30\Leftrightarrow  A\ge \dfrac{30}{(1+6{,}5\%)^3-1} \approx 144{,}26$ triệu.
	} 
\end{ex}

% nguồn: Ok_TuanLuan_Toan11_Lan2
\begin{ex}%[1D6H4-6]%[ma:9045ba20]
	Với nước biển có nồng độ muối $30\%$, nhiệt độ $T$ ($^\circ$C) của nước biển được tính bởi công thức $T = 7{,}9\ln (1{,}0245 - d) + 61{,}84$ trong đó $d$ (g/cm$^3$) là khối lượng  riêng của nước biển.
	\begin{flushright}
		(Ron Larson, Intermediate  Algebra, Cengage)
	\end{flushright}
	Biết vùng biển khơi mặt ở một khu vực có nồng độ muối $30\%$ và nhiệt độ là  $8^\circ$ C. Tính khối lượng riêng của nước biển ở vùng biển đó (làm tròn đến $4$ chữ số thập phân sau dấu phẩy).
	\loigiai{ 
		Vì vùng biển khơi mặt ở một khu vực có nồng độ muối $30\%$ và nhiệt độ là $8^\circ$ C nên $T =  8$.\\
		Thay $T = 8$ vào công thức $T = 7{,}9\ln (1{,}0245 - d)  + 61{,}84$, ta được
		\begin{eqnarray*}
			& &T = 7{,}9\ln  (1{,}0245 - d) + 61{,}84\\
			&\Leftrightarrow& 8  = 7{,}9\ln (1{,}0245 - d) + 61{,}84\\
			&\Leftrightarrow& 7{,}9\ln (1{,}0245 - d) =  -53{,}84\\
			&\Leftrightarrow& \ln (1{,}0245 - d) =  -\dfrac{2692}{395}\\
			&\Leftrightarrow& 1{,}0245 -  d = \mathrm{e}^{-\frac{2692}{395}}\\
			 &\Leftrightarrow& d = 1{,}0245 - \mathrm{e}^{-\frac{2692}{395}}\\
			&\Leftrightarrow& d  \approx 1{,}0234.
		\end{eqnarray*}
		Vậy khối lượng riêng của  nước biển ở vùng biển đó là $1{,}0234$ g/cm$^3$.
	 }
\end{ex}

% nguồn: Ok_TuanLuan_Toan11_Lan2
\begin{ex}%[1D6H4-6]%[ma:18d3b2fc]
	Tại thời điểm ban đầu nếu đầu tư $P$ đô la  với tỷ lệ lãi suất được tính gộp liên tục hàng năm không đổi là $r$ thì giá trị tương lai của khoản đầu tư này sau $t$ năm là $B(t)=P\mathrm{e}^{rt}$ đô la. Giả sử tỷ lệ lãi suất tính gộp hàng năm là $8\%$. Hỏi sau bao nhiêu năm thì số tiền đầu tư ban đầu tăng thêm ít nhất $50\%$?
	 \loigiai{
		Theo đề ra ta  có
		$P\mathrm{e}^{0{,}08t} > 1{,}5P \Leftrightarrow e^{0{,}08t}>1{,}5\Leftrightarrow 0{,}08t > \ln1{,}5\Leftrightarrow t  > \dfrac{\ln 1{,}5}{0{,}08}=5{,}06$.\\
		Vậy sau $6$ năm thì số tiền đầu tư ban đầu tăng thêm ít nhất $50\%$. 
	} 
\end{ex}

% nguồn: Ok_TuanLuan_Toan11_Lan2
\begin{ex}%[1D6H4-6]%[ma:fcb260e3]
	Cường độ ánh sáng tại độ sâu $h$ (m) dưới một mặt hồ được tính bởi công thức $ I=I_0\cdot2^{-\tfrac{h}{4}}$, trong đó $I_0$ là cường độ ánh sáng  tại mặt hồ. Tại độ sâu $3$ (m), cường độ ánh sáng giảm bao nhiêu phần trăm so với cường độ ánh sáng tại độ sâu $ 1$ (m)?
	 \loigiai{
		Cường độ ánh sáng  tại độ sâu $1$ (m)là $I_1=I_0\cdot2^{-\tfrac{1}{4}}$.\\
		Cường độ ánh sáng tại độ sâu $3$ (m)  là $I_2=I_0\cdot2^{-\tfrac{3}{4}}$.\\
		Ta có  $\dfrac{I_2}{I_1}=\dfrac{I_0\cdot2^{-\tfrac{3}{4}}}{I_0\cdot2^{-\tfrac{1}{4}}}=2^{-\tfrac{1}{2}}\approx 0{,}71=71\%$.\\
		Phần  trăm giảm cường độ ánh sáng được tính bằng công thức $\left(1-\dfrac{I_2}{I_1}\right) \cdot 100 \%$.
		Phần trăm giảm  là
		$$ 
		\left(1-\frac{1}{\sqrt{2}}\right) \cdot  100 \% \approx(1-0{,}7071) \cdot 100 \%=0{,}2929 \cdot 100 \%=29{,}29 \%
		 $$
	Vậy cường độ ánh sáng tại độ sâu $3$ (m) giảm khoảng $29$\% so với cường độ ánh sáng tại độ sâu  $1$ (m).}
\end{ex}

% nguồn: Ok_TuanLuan_Toan11_Lan2
\begin{ex}%[1D6H4-6]%[ma:907f34ea]
	Một vi khuẩn có khối lượng khoảng $5\cdot 10^{-13}$ gam và cứ sau $20$ phút vi khẩn đó tự nhân đôi một  lần. Giả sử được nuôi trong các điều kiện sinh trưởng tối ưu và mỗi con vi khuẩn đều tồn tại ít nhất là $60$ giờ. Hỏi sau bao nhiêu giờ khối lượng do tế bào vi khuẩn sinh ra sẽ đạt tới khối lượng của Trái Đất là $6\cdot 10^{27}$ gam (làm tròn kết quả đến hàng đơn vị)?
	 \loigiai{
		Gọi $\alpha_0=5\cdot 10^{-3}$ (gam) là khối  lượng ban đầu của vi khuẩn.\\
		Sau $20$ phút đầu tiên, khối lượng của vi khuẩn là $\alpha_0\cdot  2$.\\
		Sau $20$ phút thứ $2$, khối lượng của vi khuẩn  là $(\alpha_0\cdot 2)\cdot 2=\alpha_0\cdot 2^2$.\\
		Sau $20$ phút thứ  $3$, khối lượng của vi khuẩn là $(\alpha_0\cdot 2^2)\cdot 2=\alpha_0\cdot 2^3$.\\
		Sau $20$ phút thứ $n$, khối lượng vi khuẩn là  $a\cdot 2^n$.\\
		Giả sử \begin{eqnarray*}
		\alpha_0\cdot 2^n=6\cdot  10^{27}&\Leftrightarrow &5\cdot 10^{-3}\cdot 2^n=6\cdot 10^{27}\\
		&\Leftrightarrow &2^n=1{,}2\cdot 10^{30}\\ 
		&\Leftrightarrow &n=\log_2\left(1{,}2\cdot 10^{10}\right)\approx 100. 
	\end{eqnarray*}
	Vậy sau khoảng $100\cdot 20=2\,000$ phút $\approx 33{,}3$ (giờ) thì khối lượng của tế bào vi khuẩn này  sinh ra sẽ đạt tới khối lượng của Trái Đất.
 }
\end{ex}

% nguồn: Ok_TuanLuan_Toan11_Lan2
\begin{ex}%[1D6H4-6]%[ma:d8880af2]
	Một người gởi $60$ triệu đồng vào ngân hàng với lãi suất  $6\%$ một năm. Biết rằng nếu không rút tiền ra khỏi ngân hàng thì cứ sau mỗi năm số tiền lãi sẽ nhập vào gốc để tính lãi cho năm tiếp theo. Hỏi sau ít nhất bao nhiêu năm người đó nhận được số tiền nhiều hơn $100$ triệu đồng gồm cả gốc lẫn lãi?
	 \loigiai{
		Ta có $S=A(1+r)^n$.\\ 
		Để số tiền cả gốc lẫn lãi  lớn hơn $100$ triệu thì
		$$60(1+6\%)^n>100\Leftrightarrow (1+6\%)^n>\dfrac{100}{60}\Leftrightarrow  n>\log_{1+6\%}\dfrac{100}{60}\approx 8{,}766.$$
		Vậy sau ít nhất $9$ năm  thì người đó nhận được số tiền nhiều hơn $100$ triệu.
	} 
\end{ex}

% nguồn: Ok_TuanLuan_Toan11_Lan2
\begin{ex}%[1D6H4-6]%[ma:ad54bd5b]
	Bác An gửi ngân hàng số tiền $200$ triệu đồng theo thể thức lãi kép với kỳ hạn $ 6$ tháng với lãi suất $3{,}5\%/$kỳ. Số tiền cả vốn và lãi được ngân  hàng tính theo công thức $T=T_0\left(1+r\right)^n$, trong đó $T_0$ là số tiền gốc và $n$ là số kỳ đã gửi. Hỏi sau 3 năm bác An mới rút tiền thì bác thu được số tiền lãi là bao nhiêu triệu đồng?
	 \loigiai{
		Số tiền cả vốn và lãi sau  $3$ năm gửi là $T=200\left(1+0{,}035\right)^6\approx 245{,}85$ (triệu đồng).\\
	Vậy số tiền lãi mà bác An thu được sau 3 năm gửi xấp  xỉ là $ 245{,}85-200\approx 45{,}85$ (triệu đồng).}
\end{ex}

% nguồn: Ok_TuanLuan_Toan11_Lan1
\begin{bt}%[1D6V4-6]%[ma:22f5edb1]
	Sự tăng trưởng của một quần thể vi khuẩn được tính theo công thức $S=a\cdot 4^{rt}$, trong đó $a$ là số lượng  vi khuẩn ban đầu, $r$ là tỉ lệ tăng trưởng, $t$ (tính theo giờ) là thời gian tăng trưởng. Biết số lượng vi khuẩn ban đầu là $20$ con, sau $15$ giờ $320$ con. Hỏi sau bao nhiêu giờ thì số lượng vi khuẩn là $81\,920$ con?
	 \loigiai{
		Theo bài ra ta có $320=20\cdot 4^{r\cdot15}\Leftrightarrow4^{15\cdot  r}=16\Leftrightarrow 15\cdot r=2\Leftrightarrow r=\dfrac{2}{15}$.\\
		 Gọi $t_0$ là thời gian mà vi khuẩn đạt $81\,920$ con. Ta có
		$$81\,920=20\cdot 4^{\tfrac{2}{15}\cdot t_0}\Leftrightarrow 4^{\tfrac{2}{15}\cdot t_0} =4096\Leftrightarrow  \dfrac{2}{15}\cdot t_0=6\Leftrightarrow t_0=45 \text{ (giờ)}.$$
	 }
\end{bt}

% nguồn: Ok_TuanLuan_Toan11_Lan1
\begin{bt}%[1D6V4-6]%[ma:a03c7272]
	Người ta dùng thuốc để  khử khuẩn cho một thùng nước. Biết rằng nếu lúc đầu mỗi mi-li-lit nước chứa $P_{0}$ vi khuẩn thì sau $t$ giờ (kể từ khi cho thuốc vào thùng), số lượng vi khuẩn trong mỗi mi-li-lit nước còn lại là $P=P_{0} \cdot 10^{-\alpha t}, \,\alpha$ là hằng số. Biết rằng số lượng vi khuẩn trong mỗi mi-li-lit nước có $9\,000$ vi khuẩn và sau 2 giờ, số lượng vi khuẩn trong mỗi mi-li-lit nước còn lại là $6\,000$. Hỏi sau thời gian bao lâu thì số lượng vi khuẩn trong mỗi mi-li-lit nước trong thùng ít hơn hoặc bằng $1\,000$?
	\loigiai{ 
		Theo  bài ra ta có
		$6\,000=9\,000 \cdot 10^{-2 \alpha}  \Leftrightarrow 10^{-2\alpha}=\dfrac{2}{3} \Leftrightarrow \left( 10^{-\alpha} \right)^{2}=\dfrac{2}{3} \Leftrightarrow 10^{-\alpha}=\sqrt{\dfrac{2}{3}}$.\\
		Số lượng vi khuẩn trong mỗi mi-li-lit nước trong thùng ít hơn hoặc bằng $1\,000$, ta có bất phương trình 
		\begin{align*}
			 &9\,000\cdot 10^{-\alpha t} \leq 1\,000 \\
			\Leftrightarrow  &(10^{-\alpha})^{t} \leq \dfrac{1}{9}\\
			\Leftrightarrow &  \left( \sqrt{\dfrac{2}{3}} \right)^{t} \leq \dfrac{1}{9} \\
			\Leftrightarrow& t  \geq \log_{\sqrt{\dfrac{2}{3}}} \left(\dfrac{1}{9}\right) \approx 10{,}84 \,\,\text{(giờ).}
		\end{align*}
		Vậy sau khoảng hơn $10$ giờ $50$ phút thì số lượng vi khuẩn trong mỗi mi-li-lit nước trong thùng ít  hơn hoặc bằng $1\,000$.
	} 
\end{bt}

% nguồn: Ok_TuanLuan_Toan11_Lan2
\begin{ex}%[1D6V4-6]%[ma:b91b778b]
	Sự tăng trưởng của một quần thể  vi khuẩn được tính theo công thức $S=a\cdot 4^{rt}$, trong đó $a$ là số lượng vi khuẩn ban đầu, $r$ là tỉ lệ tăng trưởng, $t$ (tính theo giờ) là thời gian tăng trưởng. Biết số lượng vi khuẩn ban đầu là $20$ con, sau $15$ giờ $320$ con. Hỏi sau bao nhiêu giờ thì số lượng vi khuẩn là $81\,920$ con?
	 \loigiai{
		Theo  bài ra ta có $320=20\cdot 4^{r\cdot15}\Leftrightarrow4^{15\cdot r}=16\Leftrightarrow 15\cdot r=2\Leftrightarrow r=\dfrac{2}{15}$.\\
		Gọi $t_0$ là thời gian mà vi  khuẩn đạt $81\,920$ con. Ta có
		$$81\,920=20\cdot 4^{\tfrac{2}{15}\cdot  t_0}\Leftrightarrow 4^{\tfrac{2}{15}\cdot t_0} =4096\Leftrightarrow \dfrac{2}{15}\cdot t_0=6\Leftrightarrow t_0=45 \text{ (giờ)}.$$
	} 
\end{ex}

% nguồn: Ok_TuanLuan_Toan11_Lan2
\begin{ex}%[1D6V4-6]%[ma:e1aaf34f]
	Người ta dùng thuốc để khử khuẩn cho một thùng nước. Biết rằng nếu lúc đầu mỗi mi-li-lit nước chứa $P_{0}$ vi khuẩn thì sau $t$ giờ (kể từ khi cho thuốc vào thùng), số lượng vi khuẩn trong mỗi mi-li-lit nước còn lại  là $P=P_{0} \cdot 10^{-\alpha t}, \,\alpha$ là hằng số. Biết rằng số lượng vi khuẩn trong mỗi mi-li-lit nước có $9\,000$ vi khuẩn và sau 2 giờ, số lượng vi khuẩn trong mỗi mi-li-lit nước còn lại là $6\,000$. Hỏi sau thời gian bao lâu thì số lượng vi khuẩn trong mỗi mi-li-lit nước trong thùng ít hơn hoặc bằng $1\,000$?
	 \loigiai{
		Theo bài ra ta có 
		$6\,000=9\,000 \cdot 10^{-2 \alpha} \Leftrightarrow 10^{-2\alpha}=\dfrac{2}{3}  \Leftrightarrow \left( 10^{-\alpha} \right)^{2}=\dfrac{2}{3} \Leftrightarrow 10^{-\alpha}=\sqrt{\dfrac{2}{3}}$.\\
		Số lượng vi khuẩn trong mỗi mi-li-lit  nước trong thùng ít hơn hoặc bằng $1\,000$, ta có bất phương trình
		\begin{align*}
			&9\,000\cdot 10^{-\alpha t} \leq  1\,000 \\
			\Leftrightarrow &(10^{-\alpha})^{t} \leq \dfrac{1}{9}\\ 
			\Leftrightarrow & \left( \sqrt{\dfrac{2}{3}}  \right)^{t} \leq \dfrac{1}{9} \\
			\Leftrightarrow& t \geq \log_{\sqrt{\dfrac{2}{3}}} \left(\dfrac{1}{9}\right) \approx  10{,}84 \,\,\text{(giờ).}
		\end{align*}
		Vậy sau khoảng hơn $10$ giờ $50$ phút thì số lượng  vi khuẩn trong mỗi mi-li-lit nước trong thùng ít hơn hoặc bằng $1\,000$.
	 }
\end{ex}

% nguồn: Ok_TuanLuan_Toan11_Lan1
\begin{bt}%[1D6C4-6]%[ma:b6bc273e]
	Ngày 1/1/2024, ông Thành gửi vào ngân hàng $200$ triệu đồng theo hình thức lãi kép với lãi suất $1{,}7\% / 1$ năm, kỳ hạn $1$ tháng. Bắt đầu từ tháng $2/2024$, đều  đặn đầu mỗi tháng, ông Thành rút $4$ triệu đồng để chi tiêu. Gọi $K_n$ (triệu đồng) là số tiền còn lại trong ngân hàng sau lần rút thứ $n$ của ông Thành (mỗi lần, ông Thành rút đúng $4$ triệu đồng), giả sử lãi suất ngân hàng không thay đổi hàng năm, tìm số $n$ nhỏ nhất sao cho $K_n < 4$ (triệu đồng).
	 \loigiai{
		\begin{itemize}
			\item Số tiền còn lại trong ngân hàng  sau lần rút thứ nhất là $200\cdot \left(1+1{,}7\%\right)-4$ (triệu đồng).
			 \item Số tiền còn lại trong ngân hàng sau lần rút thứ hai là \[200\cdot \left(1+1{,}7\%\right)^2-4\cdot\left[\left(1+1{,}7\%\right)+1\right]\text{ (triệu đồng).}\]
			\item Số tiền còn lại trong ngân hàng sau lần  rút thứ ba là \[200\cdot \left(1+1{,}7\%\right)^3-4\cdot\left[\left(1+1{,}7\%\right)^2+\left(1+1{,}7\%\right)+1\right]\text{ (triệu đồng).}\]
			\item  $\dotfill$
			\item Số tiền còn lại trong ngân hàng sau lần  rút thứ $n$ là \[200\cdot \left(1+1{,}7\%\right)^n-4\cdot\left[\left(1+1{,}7\%\right)^{n-1}+\cdots+\left(1+1{,}7\%\right)+1\right]\text{ (triệu đồng).}\]
		\end{itemize}
		Để $K_n<4$ thì 
		\begin{eqnarray*}
			 & &200\cdot \left(1+1{,}7\%\right)^n-4\cdot\left[\left(1+1{,}7\%\right)^{n-1}+\cdots+\left(1+1{,}7\%\right)+1\right]<4\\
			&\Leftrightarrow& 200\cdot \left(1+1{,}7\%\right)^n-4\cdot\dfrac{(1+1{,}7\%)^n-1}{1{,}7\%}<4\\ 
			&\Leftrightarrow& n<\log_{\left(1+1{,}7\%\right)}{\dfrac{4\left(1{,}7\%-1\right)}{200\cdot 1{,}7\%-1}}\\ 
			&\Leftrightarrow&  n<111{,}5.
		\end{eqnarray*}
		 Vậy $n=9$ năm.
	 }
\end{bt}

% nguồn: Ok_TuanLuan_Toan11_Lan2
\begin{ex}%[1D6C4-6]%[ma:41b34014]
	Ngày 1/1/2024, ông Thành gửi vào ngân hàng $200$ triệu đồng theo hình thức lãi kép với lãi suất $1{,}7\% / 1$ năm, kỳ hạn $1$ tháng. Bắt đầu từ tháng $2/2024$, đều đặn đầu mỗi tháng, ông Thành rút $4$ triệu đồng để chi tiêu. Gọi $K_n$ (triệu đồng) là số tiền còn lại trong ngân hàng sau lần rút thứ $n$ của ông Thành (mỗi lần, ông Thành rút đúng $4$ triệu đồng),  giả sử lãi suất ngân hàng không thay đổi hàng năm, tìm số $n$ nhỏ nhất sao cho $K_n < 4$ (triệu đồng).
	 \loigiai{
		\begin{itemize}
			\item Số tiền còn lại trong ngân hàng sau lần rút thứ nhất là $200\cdot \left(1+1{,}7\%\right)-4$ (triệu đồng). 
			\item Số tiền còn lại trong ngân hàng sau lần rút  thứ hai là \[200\cdot \left(1+1{,}7\%\right)^2-4\cdot\left[\left(1+1{,}7\%\right)+1\right]\text{ (triệu đồng).}\]
			\item Số tiền còn lại  trong ngân hàng sau lần rút thứ ba là \[200\cdot \left(1+1{,}7\%\right)^3-4\cdot\left[\left(1+1{,}7\%\right)^2+\left(1+1{,}7\%\right)+1\right]\text{ (triệu đồng).}\]
			 \item $\dotfill$
			\item Số  tiền còn lại trong ngân hàng sau lần rút thứ $n$ là \[200\cdot \left(1+1{,}7\%\right)^n-4\cdot\left[\left(1+1{,}7\%\right)^{n-1}+\cdots+\left(1+1{,}7\%\right)+1\right]\text{ (triệu đồng).}\]
		\end{itemize}
		Để $K_n<4$  thì
		\begin{eqnarray*}
			 & &200\cdot \left(1+1{,}7\%\right)^n-4\cdot\left[\left(1+1{,}7\%\right)^{n-1}+\cdots+\left(1+1{,}7\%\right)+1\right]<4\\
			 &\Leftrightarrow& 200\cdot \left(1+1{,}7\%\right)^n-4\cdot\dfrac{(1+1{,}7\%)^n-1}{1{,}7\%}<4\\
			 &\Leftrightarrow& n<\log_{\left(1+1{,}7\%\right)}{\dfrac{4\left(1{,}7\%-1\right)}{200\cdot 1{,}7\%-1}}\\
			&\Leftrightarrow& n<111{,}5. 
		\end{eqnarray*}
		Vậy $n=9$ năm. 
	} 
\end{ex}
